% =============================================================================
%  coupledFoam - paper
%  Build: make   (latexmk -pdf paper.tex)
%  Numbers, tables and figures are generated by bench/make_report.py into
%  numbers.tex, tables/*.tex and figures/*.pdf. Every generated item is
%  guarded, so the paper compiles at every stage of the project and shows a
%  boxed "results pending" placeholder for anything that is not there yet.
% =============================================================================
\documentclass[11pt,a4paper]{article}

\usepackage[T1]{fontenc}
\usepackage[utf8]{inputenc}
\IfFileExists{lmodern.sty}{\usepackage{lmodern}}{}
\usepackage[margin=25mm]{geometry}
\setlength{\emergencystretch}{3em}
\usepackage{amsmath}
\usepackage{amssymb}
\usepackage{graphicx}
\usepackage{float}   % [H] placement in the history appendix
\usepackage{booktabs}
\usepackage{longtable}   % generated appendix: missing or stale results
\usepackage[numbers,sort&compress]{natbib}
\IfFileExists{siunitx.sty}%
  {\usepackage{siunitx}}%
  {\providecommand{\num}[1]{#1}\providecommand{\SI}[2]{#1\,#2}}
\usepackage{xcolor}
\definecolor{cflink}{rgb}{0.10,0.10,0.45}
\usepackage[colorlinks=true,linkcolor=cflink,citecolor=cflink,
            urlcolor=cflink]{hyperref}

% ---------------------------------------------------------------------------
%  Generated numbers and helpers
% ---------------------------------------------------------------------------
\IfFileExists{numbers.tex}{% Generated by bench/make_report.py - do not edit
\newcommand{\cfCommit}{8ec5dbb}
\newcommand{\cfTZeroReOneZeroZeroIterations}{58}
\newcommand{\cfTZeroReOneZeroZeroLTwou}{2.9e-06}
\newcommand{\cfTZeroReOneZeroZeroLTwov}{5.5e-06}
\newcommand{\cfTZeroReOneZeroZeroNativeIterations}{1932}
\newcommand{\cfTestTZeroReOneZeroZeroNpOnePass}{yes}
\newcommand{\cfTestTZeroReOneZeroZeroZeroNpOnePass}{yes}
\newcommand{\cfTestTestBlockFourOpsPass}{yes}
\newcommand{\cfTestTestBlockGAMGNpFourPass}{yes}
\newcommand{\cfTestTestBlockGAMGNpOnePass}{yes}
\newcommand{\cfTestTestBlockGAMGPass}{yes}
\newcommand{\cfTestTestBlockMatrixNpFourPass}{yes}
\newcommand{\cfTestTestBlockMatrixNpOnePass}{yes}
\newcommand{\cfTestTestBlockMatrixPass}{yes}
\newcommand{\cfTestTestDoubleReduceNpFourPass}{yes}
\newcommand{\cfTestTestDoubleReduceNpOnePass}{yes}
\newcommand{\cfTestTestDoubleReducePass}{yes}
\newcommand{\cfTestTestEnvPass}{yes}
\newcommand{\cfTestTestPrecisionPass}{yes}
}{}

% \cfnum{Name}: value of the generated macro \cfName, or a pending mark
\newcommand{\cfpendingmark}{\fbox{\scriptsize pending}}
\newcommand{\cfnum}[1]{%
  \ifcsname cf#1\endcsname\csname cf#1\endcsname\else\cfpendingmark\fi}
% \cfsci{Name}: as \cfnum, typeset through siunitx \num (e.g. 2.9e-06)
\makeatletter
\def\cf@na{n/a}
\newcommand{\cfsci}[1]{%
  \ifcsname cf#1\endcsname
    \edef\cf@tmp{\csname cf#1\endcsname}%
    \ifx\cf@tmp\cf@na n/a\else\expandafter\num\expandafter{\cf@tmp}\fi
  \else\cfpendingmark\fi}
\makeatother

% Boxed placeholder
\newcommand{\pending}[1]{%
  \par\noindent\fbox{\parbox{\dimexpr\linewidth-2\fboxsep-2\fboxrule\relax}{%
    \centering\small\rule{0pt}{2.2ex}\textit{Results pending.} #1\par
    \rule{0pt}{1.5ex}}}\par}

% \cfpendingitem{figure|table}{caption}{label}: compact, non-floating
% placeholder of a missing generated figure or table (numbered, so every
% \ref resolves; a run of missing items stacks as a few short lines instead
% of a cascade of full-size boxes)
\makeatletter
\def\cf@figname{figure}
\newcommand{\cfpendingitem}[3]{%
  \par\smallskip\noindent\refstepcounter{#1}\label{#3}%
  \def\cf@kind{#1}%
  {\footnotesize\fbox{\scriptsize pending}~\textbf{\ifx\cf@kind\cf@figname
    Figure\else Table\fi~\csname the#1\endcsname.} #2\par}\smallskip}

% \cffigure[width]{file stem}{caption}{label}
% (figures/<stem>.pdf, else figures/<stem>.png - the ParaView
% renders -, else the compact pending line)
\newcommand{\cffigure}[4][0.95]{%
  \IfFileExists{figures/#2.pdf}%
    {\cf@figure{#1}{figures/#2.pdf}{#3}{#4}}%
    {\IfFileExists{figures/#2.png}%
      {\cf@figure{#1}{figures/#2.png}{#3}{#4}}%
      {\cfpendingitem{figure}{#3}{#4}}}}
\newcommand{\cf@figure}[4]{%
  \begin{figure}[tbp]
    \centering\includegraphics[width=#1\linewidth]{#2}%
    \caption{#3}\label{#4}
  \end{figure}}
\makeatother

% \cftable{file stem}{caption of the placeholder}{label}
% (a generated table carries its own caption and label)
\newcommand{\cftable}[3]{%
  \IfFileExists{tables/#1.tex}%
    {\input{tables/#1.tex}}%
    {\cfpendingitem{table}{#2}{#3}}}

% \cfrenderflag{Case}: caption note of a ParaView render that is kept
% although its coupledFoam run is not from the target commit (only with
% --allow-stale; written by bench/make_report.py as \cfRenderFlag<Case>)
\newcommand{\cfrenderflag}[1]{%
  \ifcsname cfRenderFlag#1\endcsname\ \csname cfRenderFlag#1\endcsname\fi}

% Simple non-floating algorithm block (no extra packages needed). Tab stops
% (\>) are used for indentation at the start of a line only.
\newcounter{algorithm}
\newenvironment{algo}[2]{%
  \par\medskip\noindent
  \begin{minipage}{\linewidth}
  \refstepcounter{algorithm}\label{#2}%
  \hrule\smallskip\noindent\textbf{Algorithm~\thealgorithm.}\ #1\par
  \smallskip\hrule\smallskip\small
  \begin{tabbing}
  \hspace*{2em}\=\hspace*{2em}\=\hspace*{2em}\=\hspace*{2em}\=\kill}
  {\end{tabbing}\vspace*{-1ex}\hrule\end{minipage}\par\medskip}

% Notation
\newcommand{\vect}[1]{\mathbf{#1}}
\newcommand{\Uv}{\vect{U}}
\newcommand{\Sf}{\vect{S}_f}
\newcommand{\dt}{\Delta t}
\newcommand{\VSMALL}{\varepsilon_{\mathrm{v}}}
\newcommand{\code}[1]{\texttt{#1}}

% ---------------------------------------------------------------------------
\title{\textbf{coupledFoam}: a block-coupled pressure--velocity solver with
  single-precision block-multigrid linear algebra for steady incompressible
  RANS in OpenFOAM}
\author{Jonas Pangerl\\[0.5ex]
  {\small with implementation assistance by Claude (Anthropic)}}
% Title-page note of the staleness guard (bench/make_report.py, TASK 6):
% \cfStaleWarning only with --allow-stale, \cfMissingNote when results of
% the target commit are missing (shown as pending)
\makeatletter
\newcommand{\cftitlenote}{%
  \@ifundefined{cfStaleWarning}%
    {\@ifundefined{cfMissingNote}{}%
      {\\[1.5ex]\parbox{0.85\linewidth}{\centering\small\cfMissingNote}}}%
    {\\[2ex]\fcolorbox{red}{white}{\parbox{0.85\linewidth}{\centering
      \color{red}\bfseries\cfStaleWarning}}}}
\makeatother
\date{Generated at commit \texttt{\cfnum{Commit}}\cftitlenote}

\begin{document}
\maketitle

\begin{abstract}
This paper describes \textbf{coupledFoam}, a steady, incompressible, Reynolds-averaged
Navier--Stokes solver for OpenFOAM v2606 in which velocity and pressure are
solved simultaneously as one $4\times4$-block system per cell, while the
turbulence model is solved segregated and unmodified. The block matrix and
all Krylov and multigrid vectors are stored in single precision; the
right-hand side of every outer iteration is the residual $b-Ax$ of the
double-precision discretisation, and every reduction is accumulated in
double. The outer iteration is therefore a mixed-precision defect
correction whose attainable accuracy is set by the double-precision
residual, not by the single-precision solve. The nonlinear iteration is
globalised by pseudo-transient continuation with local, solution-limited
time steps and an adaptive CFL number, a physicality line search, two-tier
cell sets that locally reduce the scheme order and the time step,
a sentinel with rollback, and an exact restart. The linear systems are
solved by flexible GMRES preconditioned with an additive-correction block
algebraic multigrid that reuses OpenFOAM's agglomeration and supports V-, F-,
W- and Notay--Vassilevski K-cycles, with an Eisenstat--Walker inner
tolerance. On the lid-driven cavity at $Re=100$ the solver reproduces the
centreline profiles of the native segregated solver \code{simpleFoam} to a
relative $L_2$ difference of \cfsci{TZeroReOneZeroZeroLTwou} in
\cfnum{TZeroReOneZeroZeroIterations} outer iterations. Performance against
tuned native SIMPLE and SIMPLEC is measured in wall-clock time and CPU-hours
to an identical convergence criterion, because the target is the fastest
solver, not the one with the fewest iterations.
\end{abstract}

\tableofcontents

% =============================================================================
\section{Introduction}\label{sec:intro}
% =============================================================================

Most finite-volume solvers for incompressible flow in production use,
including the standard OpenFOAM solvers, are \emph{segregated}: the momentum
equations are solved with a lagged pressure, and a pressure (correction)
equation derived from continuity restores mass conservation. The SIMPLE
algorithm of Patankar and Spalding~\cite{Patankar1972} and its consistent
variant SIMPLEC of Van Doormaal and Raithby~\cite{VanDoormaal1984} are the
prototypes. Their cost per iteration is low and their memory footprint is
small, but the pressure--velocity coupling is only approximated by the
splitting, so the number of outer iterations grows with mesh size and with
the strength of the coupling, and under-relaxation factors have to be
tuned for every class of flow.

\emph{Coupled} solvers discretise momentum and continuity into one linear
system and solve for velocity and pressure simultaneously. The block-implicit
multigrid method of Vanka~\cite{Vanka1986} is an early example on structured
grids. Darwish, Sraj and Moukalled~\cite{Darwish2007,Darwish2009} formulated a
coupled pressure-based finite-volume solver with Rhie--Chow
interpolation~\cite{RhieChow1983} on structured and unstructured grids.
Mangani, Buchmayr and Darwish brought this approach to OpenFOAM for steady
incompressible turbulent flow~\cite{Mangani2014a}, extended it to rotating
reference frames~\cite{Mangani2014b} and to transient flows in an ALE
formulation with implicit boundary treatment~\cite{Mangani2017}; their block
system is solved with an additive-correction algebraic multigrid in the
spirit of Hutchinson and Raithby~\cite{Hutchinson1986} and Weiss et
al.~\cite{Weiss1999}. Uroi\'c and Jasak developed a block-selective algebraic
multigrid (SAMG) for the implicitly coupled pressure--velocity system in
foam-extend~\cite{Uroic2018} and its parallelisation~\cite{Uroic2021}.
ICSFoam~\cite{Oliani2023} is an OpenFOAM library for implicitly coupled,
primarily density-based simulations of high-speed flows. The price of the
coupled formulation is memory: a $4\times4$ block per cell and per face
instead of scalar matrices, and a linear system that is indefinite and
non-symmetric.

A coupled solver for steady problems is naturally combined with
pseudo-transient continuation (PTC), in which a pseudo-time derivative
regularises the Jacobian and is removed as the solution converges. Kelley and
Keyes~\cite{Kelley1998} analysed the convergence of PTC; Ceze and
Fidkowski~\cite{Ceze2011,Ceze2013AIAA,Ceze2015} developed robust CFL
strategies (SER, EXP, RDM and a modified RDM) and a physicality-based line
search for implicit high-order RANS solvers. Mulder and Van Leer's switched
evolution relaxation~\cite{MulderVanLeer1985} is the classical CFL
strategy. Inexact solves are controlled by forcing terms in the sense of
Eisenstat and Walker~\cite{EisenstatWalker1996,Knoll2004}.

Single precision halves the memory traffic of sparse linear algebra, which is
the bottleneck of memory-bound solvers, but it cannot on its own deliver the
residual levels required for converged steady solutions. Mixed-precision
iterative refinement resolves this: the correction equation is solved in low
precision while the residual is evaluated in high precision, and the limiting
accuracy is that of the residual~\cite{Baboulin2009,Carson2018}. The
increment formulation of a nonlinear outer iteration has exactly this
structure.

\paragraph{Contributions.} This paper documents coupledFoam as implemented;
every deviation from its specification is recorded in the project file
\code{DECISIONS.md} (items D-001 to D-032, referenced below). The
contributions are:
\begin{enumerate}
\item A clean-room block-coupled pressure--velocity solver for OpenFOAM v2606
  (Section~\ref{sec:disc}) in which all boundary conditions are made implicit
  through OpenFOAM's own patch linearisation, the continuity row and the
  post-solve flux are evaluated by the same face formula, and implicit MRF
  is extracted from the native MRF operator.
\item A precision design (Sections~\ref{sec:increment} and~\ref{sec:la})
  with float block storage and double residuals and reductions on a
  double-precision OpenFOAM build, together with the observations that made
  it work in practice: a strong pressure reference for closed domains, a
  best-iterate Krylov return and the treatment of failed linear solves.
\item A globalisation layer (Section~\ref{sec:algo}): PTC with a local
  solution-limited time step, mRDM/EXP/SER CFL control, physicality line
  search, static, dynamic and zonal remediation sets, sentinel and rollback,
  exact restart, Eisenstat--Walker forcing and optional Anderson acceleration.
\item A block additive-correction AMG (Section~\ref{sec:gamg}) on native
  OpenFOAM agglomeration with V-, F-, W- and K-cycles, a coarsening-ratio
  rule derived from the cycle cost, a dense coarsest-level LU, processor
  agglomeration and a cycle-efficiency controller.
\item A verification and benchmark protocol (Sections~\ref{sec:vv}
  and~\ref{sec:perf}) that compares against tuned native solvers in
  wall-clock time and CPU-hours to an identical convergence criterion.
\end{enumerate}

% =============================================================================
\section{Discretisation and the coupled system}\label{sec:disc}
% =============================================================================

\subsection{Governing equations and notation}

The solver treats the steady incompressible RANS equations for the velocity $\Uv$ and
the kinematic pressure $p$,
\begin{align}
  \nabla\cdot(\Uv_{\mathrm{r}}\Uv)
  - \nabla\cdot\bigl(\nu_{\mathrm{eff}}\nabla\Uv\bigr)
  - \nabla\cdot\bigl(\nu_{\mathrm{eff}}\,\operatorname{dev2}(\nabla\Uv^{T})\bigr)
  + \nabla p + \vect{\Omega}\times\Uv &= 0, \label{eq:mom}\\
  \nabla\cdot\Uv &= 0, \label{eq:cont}
\end{align}
with $\nu_{\mathrm{eff}}=\nu+\nu_t$, $\operatorname{dev2}(\mathsf{T}) =
\mathsf{T}-\tfrac{2}{3}\operatorname{tr}(\mathsf{T})\mathsf{I}$ as in
OpenFOAM's linear viscous stress, and, inside multiple-reference-frame (MRF)
zones rotating with $\vect\Omega$, the absolute velocity $\Uv$ convected by the
relative velocity $\Uv_{\mathrm{r}}$ (Section~\ref{sec:mrf}); outside MRF
zones $\Uv_{\mathrm{r}}=\Uv$ and $\vect\Omega=0$. The turbulence quantities
($k$, $\omega$, $\nu_t$) are frozen during the coupled solve and updated
afterwards (Section~\ref{sec:turb}).

For a cell $P$ with volume $V_P$ and face $f$ shared with neighbour $N$, the notation comprises
the face area vector $\Sf$ (pointing out of the owner), $|\Sf|$, the
interpolation weight $w_f$ of the owner, the face flux $\phi_f$
($\mathrm{m^3\,s^{-1}}$), the non-orthogonal delta coefficient
$\Delta_f = 1/(\vect{n}_f\cdot\vect{d}_f)$ of OpenFOAM, and the orientation
$s_{Pf}=+1$ if $P$ owns $f$ and $-1$ otherwise. The weight of cell $P$ on face
$f$ is $w_{Pf}=w_f$ for the owner and $1-w_f$ for the neighbour. The unknowns
per cell are $\vect{x}_P=(u,v,w,p)^T$.

\subsection{Momentum rows}

\paragraph{Convection with deferred correction.}
The implicit convection operator is OpenFOAM's \code{fvm::div} with a
first-order upwind scheme (key \code{div(phi,U)\_upwind}) on the flux of the
previous iteration: owner diagonal $+\max(\phi_f,0)$, upper coefficient
$-\max(-\phi_f,0)$, lower coefficient $-\max(\phi_f,0)$, neighbour diagonal
$+\max(-\phi_f,0)$. The user's high-order scheme (key \code{div(phi,U)},
default \code{linearUpwind grad(U)}) enters by deferred
correction~\cite{Khosla1974,Ferziger2002}:
\begin{equation}
  \vect{b}^{\mathrm{DC}}_P = -\beta_P\,V_P\Bigl[
    (\nabla\cdot(\phi\Uv))^{\mathrm{HO}}_P - (\nabla\cdot(\phi\Uv))^{\mathrm{UD}}_P
  \Bigr],\qquad \beta_P\in[0,1],
  \label{eq:dc}
\end{equation}
where both divergences are explicit evaluations (\code{fvc::div}) with the
respective scheme. The per-cell blending factor $\beta_P$ is $0$ during
start-up and in remediated cells and $1$ otherwise
(Sections~\ref{sec:startup} and~\ref{sec:remediation}); if $\beta\equiv0$ the
high-order evaluation is skipped.

\paragraph{Diffusion.}
The implicit diffusion operator is \code{fvm::laplacian} with the
\emph{uncorrected} scheme, i.e.\ the orthogonal part
$\nu_f|\Sf|\Delta_f(\Uv_N-\Uv_P)$ with $\nu_f$ the linearly interpolated
$\nu_{\mathrm{eff}}$. The non-orthogonal part is added explicitly with a
per-face limiter (Section~\ref{sec:nonorth}), and the transpose stress
$\nabla\cdot(\nu_{\mathrm{eff}}\operatorname{dev2}(\nabla\Uv^T))$ is explicit,
exactly as in the native linear viscous stress. The solver assembles the
diffusion itself instead of calling the turbulence model's
\code{divDevSigma}, because the latter takes its non-orthogonal limiter from
\code{fvSchemes} and cannot vary it per cell (D-014).

\paragraph{Pressure gradient.}
The pressure gradient is the Green--Gauss sum
$\sum_f s_{Pf}\Sf p_f$ with $p_f=w_fp_P+(1-w_f)p_N$, so face $f$ contributes
$s_{Pf}w_{Pf}\Sf$ to column~3 of the diagonal block of $P$ and
$s_{Pf}w_{Nf}\Sf$ to column~3 of the off-diagonal block $(P,N)$.

\paragraph{Pseudo-time term.}
Each momentum diagonal receives $V_P/\dt_P$ (Section~\ref{sec:ptc}); in
increment form (Section~\ref{sec:increment}) the old-time value is the
current iterate, so the term contributes only to the matrix, never to the
residual, and cannot alter the steady solution.

\subsection{Continuity row and Rhie--Chow interpolation}\label{sec:rc}

Let $a^{c}_P$, $c\in\{u,v,w\}$, be the three momentum diagonal coefficients
after all momentum contributions \emph{including} the pseudo-time term, and
\begin{equation}
  \bar a_P = \tfrac13\bigl(a^u_P+a^v_P+a^w_P\bigr),\qquad
  D_P = \frac{V_P}{\max(\bar a_P,\VSMALL)},\qquad
  D_f = w_fD_P+(1-w_f)D_N .
\end{equation}
The face flux is the Rhie--Chow flux~\cite{RhieChow1983}
\begin{equation}
  \phi_f = \Sf\cdot\bigl(w_f\Uv_P+(1-w_f)\Uv_N\bigr)
  - g_f\,(p_N-p_P) + q_f,
  \qquad g_f = D_f|\Sf|\Delta_f,
  \label{eq:phi}
\end{equation}
\begin{equation}
  q_f = D_f\Bigl[\bigl(w_f\nabla p_P+(1-w_f)\nabla p_N\bigr)\cdot\Sf
      - |\Sf|\,c_f(p)\Bigr],
  \label{eq:q}
\end{equation}
where $c_f(p)$ is the limited non-orthogonal correction of the face-normal
pressure gradient (Section~\ref{sec:nonorth}). The implicit face-gradient
term $-g_f(p_N-p_P)$ forms a Laplacian block in the pressure column of the
continuity row; the ``averaged cell gradient'' term $q_f$ is explicit. The
continuity row of cell $P$ is the discrete $\sum_f s_{Pf}\phi_f=0$ with
$\phi_f$ from~\eqref{eq:phi}, i.e.\ divergence entries
$s_{Pf}w_{Pf}\Sf^T$ (diagonal block) and $s_{Pf}w_{Nf}\Sf^T$ (off-diagonal
block) in columns 0--2, the Laplacian $+g_f$ on the diagonal and $-g_f$ off
the diagonal in column~3, and $-s_{Pf}q_f$ on the right-hand side. No
pseudo-time term is added to the continuity row.

The matrix row and the post-solve flux update are built from the same face
formula~\eqref{eq:phi} with the same $D_f$. Consequently, after the update
$\sum_f s_{Pf}\phi_f$ equals, cell by cell, the continuity residual of the
assembled row evaluated at the new state (with the explicit term $q_f$
refreshed): the mass imbalance of the output flux is exactly the continuity
part of the outer residual and vanishes with it. The dimensionally consistent
coefficient $D_f|\Sf|\Delta_f$ (rather than $D_f|\Sf|^2\Delta_f$) and the
sign convention of an M-matrix Laplacian are recorded in D-003.

\paragraph{Boundary faces of the Rhie--Chow terms.}
On physical boundary faces $D_f=D_P$ where the pressure boundary condition
fixes the value (\code{fixesValue()}, e.g.\ \code{fixedValue},
\code{totalPressure}), and $D_f=0$ elsewhere, since an averaged-gradient term
on a zero-gradient wall would create a spurious wall flux (D-013). On
symmetry, symmetry-plane and wedge patches the face flux is identically zero
and the patch contributes nothing to the continuity row. Processor faces use
the interpolated $D_f$ and are implicit; cyclic and cyclicAMI faces use the
interpolated $D_f$ with lagged neighbour values.

\subsection{Block structure}\label{sec:block}

With the notation above, the diagonal block of cell $P$ and the off-diagonal
block of an internal face $f$ between $P$ and $N$ are
\begin{equation}
  \mathsf{A}_{PP} =
  \begin{pmatrix}
    a^u_P & 0 & 0 & \gamma^x_P\\
    0 & a^v_P & 0 & \gamma^y_P\\
    0 & 0 & a^w_P & \gamma^z_P\\
    \delta^x_P & \delta^y_P & \delta^z_P & \sum_f g_f
  \end{pmatrix}
  + V_P\begin{pmatrix}[\vect\Omega]_\times & 0\\ 0 & 0\end{pmatrix},
  \qquad
  \mathsf{A}_{PN} =
  \begin{pmatrix}
    a_{PN}\mathsf{I}_3 & s_{Pf}w_{Nf}\Sf\\
    s_{Pf}w_{Nf}\Sf^T & -g_f
  \end{pmatrix},
  \label{eq:block}
\end{equation}
with $\vect\gamma_P=\sum_f s_{Pf}w_{Pf}\Sf$ plus boundary terms, $\vect\delta_P$
the corresponding divergence row (equal to $\vect\gamma_P$ on internal faces),
$a_{PN}$ the upwind-plus-diffusion coupling, and $[\vect\Omega]_\times$ the
cross-product matrix. The pressure-gradient column of the momentum rows and
the velocity-divergence row are transposes of each other; the $(3,3)$ entries
are the Rhie--Chow stabilisation of the otherwise saddle-point system. The
blocks are stored row-major, 16 values per cell and per face
(Section~\ref{sec:storage}); rows 0--2 are momentum, row~3 is continuity.

\paragraph{Continuity-row scaling.}
The continuity row is multiplied by
\begin{equation}
  s_p = \frac{1}{\max(\langle D\rangle,\VSMALL)},\qquad
  \langle D\rangle = \text{global mean of } D_P,
\end{equation}
before the system is narrowed to single precision, so that the magnitudes of
the continuity and momentum rows are comparable. Since this is a pure row
scaling, the solution vector needs no unscaling; the residual norms of
Section~\ref{sec:norms} are those of the scaled system.

\subsection{Boundary conditions}\label{sec:bc}

The specification tabulated the implicit treatment of a small set of
boundary-condition pairs, and the table was partly lost in its source
(D-002). The implementation therefore treats every physical patch uniformly
through the linearisation that every OpenFOAM patch field provides,
\begin{equation}
  \Uv_b = \vect\alpha^U_b\odot\Uv_P + \vect\beta^U_b,\qquad
  p_b = \alpha^p_b\,p_P + \beta^p_b,
  \label{eq:bclin}
\end{equation}
with $\vect\alpha,\vect\beta$ from \code{valueInternalCoeffs(w)} and
\code{valueBoundaryCoeffs(w)} (componentwise product $\odot$). The momentum
rows take the native \code{internalCoeffs}/\code{boundaryCoeffs} of the upwind
convection and uncorrected Laplacian operators, and the pressure gradient
contributes $\Sf\alpha^p_b$ to column~3 of the diagonal block and
$-\Sf\beta^p_b$ to the right-hand side. The continuity row takes
$\Sf^T\operatorname{diag}(\vect\alpha^U_b)$ in columns 0--2,
$-\Sf\cdot\vect\beta^U_b$ on the right-hand side, and on fixed-pressure faces
the Rhie--Chow term $-g_b(p_b-p_P)$ linearised with~\eqref{eq:bclin}.
Because the linearisation is native, \code{fixedValue}, \code{zeroGradient},
\code{slip}/\code{symmetry}, \code{inletOutlet}, \code{totalPressure} and any
condition derived from fixed-value, zero-gradient or mixed patches work
without special code (the mixed value fraction of \code{inletOutlet} and the
dynamic head of \code{totalPressure} are evaluated from the previous outer
iteration). Wall functions act through $\nu_t$ on the wall face. For
outflow through a fixed-value face the native convection boundary term is
kept (D-019).

Processor patches are fully implicit: their coefficients are stored as block
interface coefficients and applied during the matrix--vector product
(Section~\ref{sec:impl}). Cyclic, cyclicAMI and processor-cyclic patches
are explicitly coupled in this version: neighbour values are lagged from the
previous outer iteration and enter the right-hand side; implicit block
coupling across them is left to Phase~2.

For closed domains ($p$ needs a reference), the continuity row of the
reference cell is replaced by $d\,(p_{\mathrm{ref}}-p_{\mathrm{value}})=0$,
where $d$ is its former diagonal (D-021). This is exact, because the
continuity equations of a closed incompressible domain sum to zero and one of
them is redundant. The native weak reference, which doubles the diagonal of
the reference cell, leaves a near-null space; once the right-hand side is
rounded to single precision its compatibility condition is slightly violated
and the float Krylov solve drifts along the constant-pressure mode. This was
observed on the cavity case and removed by the strong reference.

\subsection{Implicit MRF}\label{sec:mrf}

The native MRF formulation solves for the absolute velocity, convects it with
the relative flux, and adds $\vect\Omega\times\Uv$ as an explicit source; this
equals the relative-frame Coriolis and centrifugal terms
$2\vect\Omega\times\Uv_{\mathrm{r}}+\vect\Omega\times(\vect\Omega\times\vect r)$.
coupledFoam makes the same term implicit by adding
$V_P[\vect\Omega]_\times$ to the $3\times3$ momentum part of the diagonal block
(D-006; a factor $2$ would be wrong for the absolute velocity). The per-cell
matrix is extracted from the native operator \code{MRFZoneList::DDt} applied to
the three unit vector fields, so any number of zones and time-dependent
$\vect\Omega$ are handled through the public API only; it is re-extracted every
outer iteration. As in \code{simpleFoam}, the flux is stored relative
(\code{MRF.makeRelative} after the flux update), because the turbulence model
convects with the same flux object (D-014, addendum).

\subsection{Non-orthogonal correction with a per-face limiter}\label{sec:nonorth}

For a face-normal gradient the full face value is
$\partial_n\psi|_f=\Delta_f(\psi_N-\psi_P)+c_f(\psi)$. The correction $c_f$ is
the native \code{limited corrected} $\lambda$ correction: the explicit
non-orthogonal part of the \code{corrected} scheme, limited so that its
magnitude does not exceed $\lambda/(1-\lambda)$ times the orthogonal part.
Two limiters are evaluated, the global \code{nonOrthLimiter} (default $0.5$)
and the static-remediation limiter (default $0.2$); faces adjacent to a
static-remediation cell take the latter (Section~\ref{sec:remediation}). The
same correction is used for the momentum Laplacian, where it enters as the
explicit source $\nabla\cdot(\nu_f|\Sf|c_f(\Uv))$ (zero on physical boundary
faces), and for the Rhie--Chow pressure Laplacian through $q_f$
in~\eqref{eq:q}. On an exactly orthogonal mesh, detected from the native
correction vectors, both are skipped.

\subsection{Turbulence coupling}\label{sec:turb}

The turbulence model is used unmodified through OpenFOAM's incompressible
turbulence-model interface (e.g.\ $k$-$\omega$ SST~\cite{Menter1994}, GEKO
\cite{Menter2025GEKO}); \code{turbulence->correct()} is called once per outer
iteration after the coupled update, with the native solvers and relaxation
factors for $k$ and $\omega$. Afterwards $k\ge k_{\min}$ and
$\omega\ge\omega_{\min}$ are enforced and $\nu_t$ is capped at
$10^8\,\nu$ (D-040; the number of capped cells is logged).

% =============================================================================
\section{Solution algorithm}\label{sec:algo}
% =============================================================================

\subsection{Increment formulation and the double-precision residual}\label{sec:increment}

Let $\mathcal{A}(\vect x)$ be the block operator assembled at the state
$\vect x=(\Uv,p)$ (convection linearised on the lagged flux, explicit terms
evaluated at $\vect x$) and $\vect b(\vect x)$ its right-hand side. Assembly
is in double precision. The residual
\begin{equation}
  \vect r = \mathsf{S}\bigl(\vect b-\mathcal{A}\vect x\bigr),\qquad
  \mathsf{S}=\operatorname{diag}(1,1,1,s_p)\ \text{per cell},
\end{equation}
is accumulated in double from the double coefficients \emph{before} the
coefficients are narrowed to single precision (D-012). The single-precision
matrix $\mathsf{A}_{\mathrm{f}}=\operatorname{fl}_{32}(\mathsf{S}\mathcal A)$
is used only to compute an increment,
\begin{equation}
  \mathsf{A}_{\mathrm{f}}\,\delta\vect x \approx \operatorname{fl}_{32}(\vect r),
  \qquad
  \vect x\leftarrow\vect x+\omega\,\delta\vect x,
  \label{eq:increment}
\end{equation}
with the line-search factor $\omega$ of Section~\ref{sec:ls}. This is a
mixed-precision iterative refinement~\cite{Baboulin2009,Carson2018}
superposed on the nonlinear (Picard/defect-correction) iteration: the
single-precision matrix, the inexact Krylov solve and the pseudo-time term
change only the convergence rate, while the fixed point is that of the double
residual. On the cavity test this is what allows $R<10^{-8}$ with a float
linear solver. The narrowing clamps values with $|v|>10^{30}$ to
$\pm10^{30}$ and counts them; a non-zero count is a warning and a test
failure.

Because the pseudo-time term is added with the current iterate as the old
value, it contributes $(V_P/\dt_P)\,\delta\Uv_P$ to the matrix and nothing to
$\vect r$. In increment form the row scaling and the continuity row need no
special treatment.

\subsection{Residual norms}\label{sec:norms}

With $\vect r$ and $\mathcal A\vect x$ of the scaled system, the solver uses
\begin{equation}
  \nu_{\mathrm{n}} = \sum_i\bigl(|(\mathcal A\vect x)_i|+|b_i|\bigr)+\epsilon,
  \qquad
  R_n = \frac{\|\vect r_n\|_2/\nu_{\mathrm{n},n}}{\|\vect r_1\|_2/\nu_{\mathrm{n},1}},
\end{equation}
summed over all four components of all cells, the analogue of OpenFOAM's
\code{normFactor}. $R_n$ is the combined residual that drives the CFL control
and the stop criterion (D-016). For comparison with native solvers the log
also reports the $L_1$-normalised residuals $r_U$ (rows 0--2) and $r_p$
(row~3), $\sum|r_i|/\sum(|(\mathcal A\vect x)_i|+|b_i|)$, which are
directly comparable with the ``initial residual'' of \code{simpleFoam}.

\subsection{Outer iteration}

Algorithm~\ref{alg:outer} summarises one outer iteration. Assembly is split
into two stages, because the solution-limited time step
(Section~\ref{sec:loclim}) needs the momentum residual before the pseudo-time
term is known, whereas $D=V/\bar a$ needs the pseudo-time term.

\begin{algo}{Outer iteration $n$ of coupledFoam.}{alg:outer}
$\nu_{\mathrm{eff}}$; $\beta_P$ (start-up, remediation, zonal); CFL factors $f_P$\\
\textbf{repeat} \quad \textit{(at most $1+{}$\code{maxCflCuts} times)}\\
\> $V_P/\dt_P \leftarrow \lambda_P/(\mathrm{CFL}\,f_P)$ \quad \textit{(local time step, eq.~\ref{eq:dt})}\\
\> assemble momentum rows; momentum residual $\vect r^{U}_P$ without pseudo-time term\\
\> apply the solution-limited cap to $V_P/\dt_P$ \quad \textit{(eq.~\ref{eq:loclim})}\\
\> add $V_P/\dt_P$; $\bar a_P$, $D_P$, $D_f$, $q_f$; assemble continuity row; reference cell\\
\> narrow to float; scale row 3; $\vect r=\mathsf S(\vect b-\mathcal A\vect x)$ in double; $R_n$\\
\> solve $\mathsf A_{\mathrm f}\,\delta\vect x=\vect r$ \quad \textit{(FGMRES + block GAMG, rel.\ tolerance $\eta_n$)}\\
\> \textbf{if} the solve failed ($\eta_n\|\vect r_0\|$ and absolute floor not reached in \code{maxIter}, or non-finite):\\
\> \> \textbf{if} \code{maxLinFails} consecutive failures: diagnostic exit \quad \textit{(Section~\ref{sec:sentinel})}\\
\> \textbf{else}: $\omega\leftarrow$ physicality line search \quad \textit{(eq.~\ref{eq:ls})}\\
\> \textbf{if} ($\omega<\omega_{\min}$ or failed) and cuts $<$ \code{maxCflCuts}:\\
\> \> $\mathrm{CFL}\leftarrow\max(\mathrm{CFL}_{\min},\kappa\,\mathrm{CFL})$; \textbf{repeat}\\
\textbf{if} failed: skip the update; $\mathrm{CFL}\leftarrow\max(\mathrm{CFL}_{\min},0.25\,\mathrm{CFL})$ \quad \textit{(only if \code{maxCflCuts} $<$ \code{maxLinFails})}\\
\textbf{else if} $\omega<\omega_{\min}$: $\omega\leftarrow\omega_{\min}$; mark offending cells dynamic\\
clip increments of dynamic cells; store rollback copy; $\Uv\mathrel{+}=\omega\,\delta\Uv$, $p\mathrel{+}=\omega\,\delta p$\\
update boundary values; $\phi_f$ from eq.~\ref{eq:phi}; make $\phi$ relative in MRF zones\\
sentinel check, rollback on failure \quad \textit{(Section~\ref{sec:sentinel})}\\
dynamic remediation update; \code{turbulence->correct()}; bounds; sentinel check\\
CFL strategy update with $R_n$; start-up switch; log; write; convergence test
\end{algo}

\subsection{Pseudo-transient continuation and local time step}\label{sec:ptc}

The local pseudo-time step combines a convective and a diffusive limit,
\begin{equation}
  \frac{V_P}{\dt_P} = \frac{\max(\lambda_P,\VSMALL)}{\max(\mathrm{CFL}\,f_P,\VSMALL)},
  \qquad
  \lambda_P = \tfrac12\sum_f|\phi_f| + \nu_{\mathrm{eff},P}\,V_P^{1/3},
  \label{eq:dt}
\end{equation}
with the global adaptive CFL number and the product $f_P\le1$ of the
remediation and zonal CFL factors (Section~\ref{sec:remediation}). The global
CFL follows one of three strategies~\cite{Ceze2011,Ceze2013AIAA,MulderVanLeer1985}:
\begin{align}
  \text{mRDM (default):}\quad &
  \mathrm{CFL}_{n+1} = \begin{cases}
    \min\!\bigl(\mathrm{CFL}_{\max},\ \mathrm{CFL}_n\max\bigl(1,\min(\beta_{\max},(R_{n-1}/R_n)^{\gamma})\bigr)\bigr) & R_n\le R_{n-1},\\
    \mathrm{CFL}_n & \text{otherwise},\end{cases}\\
  \text{EXP:}\quad & \mathrm{CFL}_{n+1}=\min(\mathrm{CFL}_{\max},\beta_{\exp}\mathrm{CFL}_n),\\
  \text{SER:}\quad & \mathrm{CFL}_{n+1}=\min\bigl(\mathrm{CFL}_{\max},\max(\mathrm{CFL}_{\min},\mathrm{CFL}_n R_{n-1}/R_n)\bigr),
\end{align}
with defaults $\gamma=1$, $\beta_{\max}=1.5$, $\beta_{\exp}=1.1$,
$\mathrm{CFL}_0=5$, $\mathrm{CFL}_{\min}=1$, $\mathrm{CFL}_{\max}=500$. mRDM
and EXP never decrease the CFL number; decreases come only from the line
search, failed solves and the sentinel. After every decrease a hold of
$n_{\mathrm{hold}}=10$ iterations forbids increases. With the linear-solve
failure trigger of amendment B4 the recovery after a decrease is governed by
mRDM (bounded by $\beta_{\max}$), not by a fixed ramp; after
\code{maxLinFails} (3) consecutive failed linear solves the run takes the
diagnostic exit of Section~\ref{sec:sentinel}.

\subsection{Solution-limited local time step}\label{sec:loclim}

The momentum residual of the current assembly (without the pseudo-time term,
but including the pressure gradient) estimates the local change a step would
cause, $\delta U_P = |\vect r^{U}_P|\,\dt_P/V_P$. If
$\delta U_P > f_{\mathrm{loc}}U_{\mathrm{ref}}$ ($f_{\mathrm{loc}}=0.5$), the
cell's time step alone is reduced,
\begin{equation}
  \frac{V_P}{\dt_P}\leftarrow\frac{V_P}{\dt_P}\,
  \frac{\delta U_P}{f_{\mathrm{loc}}U_{\mathrm{ref}}},
  \label{eq:loclim}
\end{equation}
so that the global CFL can stay aggressive while individual cells are
throttled; the number of limited cells is logged.

\subsection{Physicality line search}\label{sec:ls}

Given the increment, the step length is
\begin{equation}
  \omega = \min\Bigl(1,\ \min_P\frac{f_UU_{\mathrm{ref}}}{\max(|\delta\Uv_P|,\VSMALL)},\
  \min_P\frac{f_p\,p_{\mathrm{ref}}}{\max(|\delta p_P|,\VSMALL)}\Bigr),
  \label{eq:ls}
\end{equation}
reduced over all ranks, with $f_U=0.3$, $f_p=0.5$,
$U_{\mathrm{ref}}=\max|\Uv|$ (internal and physical boundary values) at the
first iteration and $p_{\mathrm{ref}}=\tfrac12U_{\mathrm{ref}}^2$, both frozen
and stored for restart~\cite{Ceze2013AIAA}. If $\omega<\omega_{\min}=0.1$, the
CFL number is cut by $\kappa=0.5$ and the step is repeated from the same
fields (reassembly with the new $\dt$), at most three times; then the step is
accepted with $\omega=\omega_{\min}$ and every cell whose increment would
have violated the bounds at $\omega_{\min}$ is marked into the dynamic
remediation set.

A linear solve that does not reach its target $\eta_n\|\vect r_{\mathrm{lin},0}\|$
(or the absolute tolerance floor) within \code{maxIter} iterations, or that
produces non-finite values, is a failed step (amendment B4, D-029, which
replaces the earlier ``no residual reduction'' rule of D-020): near the
float floor a Krylov solve can diverge, and applying even $\omega_{\min}$ of
such an increment destroys a converged state. A failed solve cuts the CFL
number by $\kappa$ and repeats the outer step from unchanged fields; the
repeat counts toward \code{maxCflCuts}. After \code{maxLinFails} (3)
consecutive failures the run leaves through the diagnostic path of
Section~\ref{sec:sentinel}. The fallback of D-020 (skip the step without a
field update and cut the CFL number by the sentinel factor) remains only for
settings with \code{maxCflCuts} $<$ \code{maxLinFails}. In addition, block
BiCGStab returns the iterate with the smallest residual if it does not
converge, so its output is never worse than the initial guess (D-020, item~1);
the GMRES variants minimise the residual over their Krylov space by
construction.

\subsection{Start-up}\label{sec:startup}

Fresh runs are initialised from \code{potentialFoam} (called by the case
script with \code{coupled.potentialInit yes}); in the benchmark its wall and
CPU time are counted as part of the coupledFoam solution
(Section~\ref{sec:perf}). For the first \code{startupUpwindIters} (50) iterations, or until
$R_n<10^{-2}$, $\beta_P=0$ everywhere, i.e.\ the discretisation is pure upwind;
afterwards $\beta_P=1$ except in remediated cells. On restart the start-up
phase is skipped.

\subsection{Remediation cell sets}\label{sec:remediation}

\paragraph{Static set.} Built once before the first iteration and split
into four categories, each with its own switch, criteria and treatment in
the case dictionary (D-066; D-063 asked for every such setting to be a
run-time keyword):
\begin{itemize}
\item \emph{meshQuality}: a face with non-orthogonality above $85^\circ$ or
  skewness above 6, a volume ratio above 30 to a neighbour, or an aspect
  ratio above 2000 (OpenFOAM's per-face quality measures; thresholds of
  D-047);
\item \emph{badMesh}: severe defects: a volume ratio above 50 to a neighbour
  (volume jump), open cells whose face area vectors do not close, and
  optionally faces beyond severe non-orthogonality or skewness limits;
\item \emph{processor}: cells with both a processor and a cyclicAMI face,
  optionally a number of cell layers next to every processor boundary;
\item \emph{wall}: optionally ``starved'' wall cells with at most two
  internal neighbours (off by default) and cell layers next to named
  patches.
\end{itemize}
The first two categories hold defective cells and receive the \emph{full}
treatment: $\beta_P=\min(\beta_P,\beta_{\mathrm{cat}})$ with
$\beta_{\mathrm{cat}}=0$ (upwind), the stronger non-orthogonal limiter $0.2$
on their faces, a limited (\code{cellLimited Gauss linear 1}) pressure
gradient in the Rhie--Chow term, and a CFL factor $0.5$. Processor and wall
cells are selected by position, not by quality, and receive a \emph{mild}
treatment by default: the CFL factor only, with the full convection scheme
and no limiters. First-order convection on a dozen well-shaped wall cells of
the motorbike body changed $C_l$ by $-18\,\%$ (D-061), which is why the
position-based categories keep the scheme. A cell in several categories
receives the most restrictive value of each parameter. The gradient inside
the user's high-order convection scheme cannot be overridden per cell; with
$\beta_{\mathrm{cat}}=0$ it does not enter (D-018). Table~\ref{tab:remcat}
lists the cells per category of every test run.

\paragraph{Dynamic set.} Updated every iteration. A cell is marked if
$|\Uv_P|>4U_{\mathrm{ref}}$, $|p_P|>15p_{\mathrm{ref}}$ (D-044),
$|\Uv_P-\bar{\Uv}_{N(P)}|>0.5\,U_{\mathrm{ref}}$ (face-neighbour mean), any of
$\Uv,p,k,\omega$ is non-finite, or it is passed in by the line search or the
sentinel. Marked cells plus one layer of face neighbours (across processor
boundaries) receive $\beta_P=0$, a CFL factor $0.1$ and increment clipping,
\begin{equation}
  |\Uv_P+\omega\,\delta\Uv_P-\bar\Uv_{N(P)}|\le c_{\mathrm{spike}}U_{\mathrm{ref}},
  \qquad c_{\mathrm{spike}}=0.5,
\end{equation}
which rescales the increment and never overwrites the field. A cell leaves
the set after 20 consecutive iterations without being re-marked. If the
union of both sets exceeds 1\,\% of the cells a warning is issued.

\paragraph{Zonal factors.} Optionally (amendment B6), multipliers on the local
time step and on $\beta_P$ are applied in named cell zones and within
$n_{\mathrm{layers}}$ face-neighbour layers of named patches (layers computed
once by a breadth-first search from the patch faces, across processor
faces); several matching entries multiply. The effective blending factor and
CFL factor are
\begin{equation}
  \beta_P = z^{\beta}_P\cdot\begin{cases}0 & P\ \text{dynamic},\\
  \min(\beta_{\mathrm{glob}},\beta_{\mathrm{static}}) & P\ \text{static},\\
  \beta_{\mathrm{glob}} & \text{otherwise},\end{cases}\qquad
  f_P = z^{\mathrm{CFL}}_P\min\bigl(1,\,f^{\mathrm{static}}_P,\,f^{\mathrm{dyn}}_P\bigr),
\end{equation}
with $\beta_{\mathrm{glob}}\in\{0,1\}$ from the start-up switch and
$f^{\mathrm{static}}_P$, $f^{\mathrm{dyn}}_P$ equal to one outside the
respective sets. Local pseudo-time is implicit local under-relaxation;
explicit per-patch relaxation factors are deliberately not provided, because
they would break the increment consistency of the line search, and processor
boundaries need no relaxation since the block interfaces are fully implicit.

\subsection{Sentinel, rollback and failure handling}\label{sec:sentinel}

Before every field update, $\Uv,p,\phi$ and (if present) $k,\omega,\nu_t$ are
copied. After the update and again after the turbulence correction a single
pass computes the extrema of $p$, $|\Uv|$, $k$, $\omega$ and the number of
non-finite values. The check fails if any value is non-finite,
$\max|\Uv|>10\,U_{\mathrm{ref}}$ or $\max|p|>50\,p_{\mathrm{ref}}$. On failure the
copies are restored, the CFL number is cut to
$\max(\mathrm{CFL}_{\min},0.25\,\mathrm{CFL})$, the offending cells are marked
into the dynamic set and a rollback is counted. After five consecutive
rollbacks the restored fields are written to a time directory
\code{<iter>\_lastValid} together with the remediation flag and a
\code{sentinelFlag} field of the offending cells, their centres are printed,
and the run stops with a fatal error. A time directory containing
non-finite values is never written.

\subsection{Restart}\label{sec:restart}

At every write time the native fields ($\Uv,p,\phi,k,\omega,\nu_t$) are written,
followed by the dictionary \code{coupledState}: iteration, CFL, previous
$R$, hold counter, start-up flag, $R_1$, $U_{\mathrm{ref}}$, $p_{\mathrm{ref}}$,
dynamic set and ages, sentinel counters, and a placeholder for an AMR
refinement history. The field $D$ of the last assembly is written alongside.
\code{coupledState} is written to a temporary file and renamed atomically,
last, so its presence marks a complete time directory (D-017). On restart
the following are restored: the fields $\Uv$, $p$, $\phi$ and the turbulence
fields; the PTC state (iteration, CFL, previous $R$, hold counter, $R_1$,
start-up flag, $U_{\mathrm{ref}}$, $p_{\mathrm{ref}}$); the dynamic set and
its ages; the sentinel counters; the Eisenstat--Walker state; the state of
the cycle-efficiency controller (current cycle type, post-smoothing count,
window and hysteresis); and the window of the force-coefficient convergence
monitor. Deliberately not restored is the Anderson history
(Section~\ref{sec:anderson}): a restarted run starts with an empty history
(D-026), so a run with Anderson acceleration is reproducible across a
restart only up to that history. $\phi$ is recomputed from the read $\Uv,p$
with the stored $D$ by~\eqref{eq:phi}; the maximum relative difference to
the read flux is logged as the flux-consistency metric (required below
$10^{-6}$ in the restart test).

\subsection{Eisenstat--Walker forcing}\label{sec:ew}

With \code{adaptiveRelTol} (default), the relative tolerance of the linear
solve in outer iteration $n$ is the forcing term of Eisenstat and Walker's
Choice~2~\cite{EisenstatWalker1996},
\begin{equation}
  \eta_n = \gamma_{\mathrm{EW}}\Bigl(\frac{R_n}{R_{n-1}}\Bigr)^{\alpha_{\mathrm{EW}}},
  \qquad \gamma_{\mathrm{EW}}=0.9,\ \alpha_{\mathrm{EW}}=2,
\end{equation}
with the safeguards, in this order: (i) if
$\gamma_{\mathrm{EW}}\eta_{n-1}^{\alpha_{\mathrm{EW}}}>0.1$ then
$\eta_n\leftarrow\max(\eta_n,\gamma_{\mathrm{EW}}\eta_{n-1}^{\alpha_{\mathrm{EW}}})$;
(ii) $\eta_n\leftarrow\min(\eta_{\max},\max(\eta_{\min},\eta_n))$
with $\eta_{\min}=10^{-3}$, $\eta_{\max}=0.5$; (iii) during start-up and
(iv) at $n=1$, $\eta_n=\eta_{\max}$. The absolute tolerance remains a floor
and at least one Krylov iteration is always performed, so the increment is
never zero. The minimum-iteration rule matters: if the absolute linear
tolerance is reached before the outer target, a zero increment freezes the
outer iteration (observed on the cavity: $R$ frozen at
\cfsci{ExpDTwoTwoFrozenR}, D-022).

\subsection{Anderson acceleration}\label{sec:anderson}

Optionally (default off), the accepted outer update is accelerated by
Anderson mixing~\cite{Anderson1965} in the Type~II form of Walker and
Ni~\cite{WalkerNi2011}. The outer iteration is read as the fixed-point map
$\vect x\mapsto G(\vect x)=\vect x+\vect f(\vect x)$ with
$\vect f=\omega\,\delta\vect x$ the accepted (line-searched, clipped)
increment. With the last $m$ differences
$\Delta\vect x_j=\vect x_{j+1}-\vect x_j$ and
$\Delta\vect f_j=\vect f_{j+1}-\vect f_j$,
\begin{equation}
  \vect\gamma=\arg\min_{\vect\gamma}\Bigl\|\mathsf W\Bigl(\vect f_k-\sum_j\gamma_j\Delta\vect f_j\Bigr)\Bigr\|_2,
  \qquad
  \vect x_{k+1}=\vect x_k+\beta\vect f_k-\sum_j\gamma_j\bigl(\Delta\vect x_j+\beta\Delta\vect f_j\bigr),
\end{equation}
where $\mathsf W$ scales velocities by $1/U_{\mathrm{ref}}$ and pressure by
$1/p_{\mathrm{ref}}$; for $\beta=1$ this is
$\vect x_{k+1}=G(\vect x_k)-\sum_j\gamma_j\Delta G_j$. The least-squares
problem is solved with a thin QR factorisation updated column by column
(classical Gram--Schmidt with one re-orthogonalisation, all dot products in
double and reduced over ranks; the oldest column is removed by Givens
rotations); a numerically dependent new column is not added. Safeguards: the
update is skipped (history kept) if $\max_j|\gamma_j|>10$; the extrapolated
state must pass the sentinel, otherwise the un-extrapolated state is
restored and the history flushed. The history is further flushed on every
CFL cut and every skipped step, when the line search returns the minimum
step $\omega_{\min}$, on every rollback, whenever the dynamic remediation set
changes (detected by a version counter of the set, not by its size), at the
start-up switch ($\beta$ from 0 to 1, which changes the discretisation), and
whenever the cycle-efficiency controller changes the preconditioner (which
changes the inner solve and hence the map). Convergence theory
for contractive and non-contractive operators and for Picard iterations of
the Navier--Stokes equations is given in~\cite{PollockRebholz2021,Pollock2019}.
The memory cost is $(2m+2)$ double 4-vectors per cell, 320\,B/cell for $m=4$
(\cfnum{BSevenAndersonGB}\,GB at 45\,M cells, Table~\ref{tab:memory}),
which is why the option is off by default and not used above about 35\,M
cells.

% =============================================================================
\section{Linear algebra}\label{sec:la}
% =============================================================================

\subsection{Storage and precision}\label{sec:storage}

The specification asked for OpenFOAM's mixed-precision build (\code{SPDP})
with single-precision \code{solveScalar}. In OpenFOAM v2606 the opposite
holds: \code{SPDP} stores fields in float and solves in double, and no build
option gives float solve with double fields (D-001). OpenFOAM is therefore
built in double precision, and the library defines its own types:
\code{blockScalar} = \code{float} for block coefficients and all Krylov and
multigrid vectors, and \code{reduceScalar} = \code{double} for every
accumulation (dot products, norms, Galerkin sums) and every MPI reduction.
The native \code{solveScalar} is never used for block data, so the library is
correct under DP, SPDP and SP builds, and the native reference solver runs
with the same double-precision fields.

The block matrix uses OpenFOAM's LDU addressing: 16 floats per cell for the
diagonal blocks (row-major, index $16c+4r+j$), 16 per internal face for the
upper (owner row, neighbour column) and 16 for the lower blocks, 4 per cell
for vectors, and 16 per face for each processor interface. On a hexahedral
mesh with about three internal faces per cell this is about
$64+2\cdot3\cdot64=448$\,B per cell, half of the double-precision
equivalent. Table~\ref{tab:memory} gives the budget for the largest target
mesh.

\begin{table}[tbp]
\centering\small
\caption{Memory budget at 45\,M cells, worst case (amendment B7). The
45\,M-cell run itself is a user acceptance step outside the one-shot scope;
Figure~\ref{fig:memory} compares the measured peak RSS of the benchmark
cases with this budget scaled per cell. The Anderson row uses the
implemented 320\,B/cell (D-026); the 11.5\,GB of the B7 table assumed
256\,B/cell.}
\label{tab:memory}
\begin{tabular}{lr}
\toprule
item & approx.\ GB\\
\midrule
mesh, fields, addressing (double) & 40--45\\
block matrix $4\times4$ (single) & 20\\
block-GAMG hierarchy, $C_{\mathrm{op}}\le1.5$ (single) & 20--30\\
FGMRES, restart 6: 12 block vectors (single) & 8.6\\
K-cycle workspace, all levels & 5\\
turbulence, rollback copies, reserve & 10\\
\midrule
total & 105--120\\
Anderson, $m=4$ (optional) & $+$\cfnum{BSevenAndersonGB}\\
\bottomrule
\end{tabular}
\end{table}

\paragraph{Double-precision reductions.}
All inner products and norms go through a small \code{doubleReduce} layer
that accumulates float products in double locally and reduces with
\code{MPI\_DOUBLE}; the multi-value variants reduce several sums in one call
to save latency. A float sum over $10^7$--$10^8$ entries loses three to four
digits and $\|\cdot\|^2$ can overflow. The unit test sums $10^8$ copies of
$\operatorname{fl}_{32}(10^{-4})$: the double accumulator reproduces
$10^8\cdot\operatorname{fl}_{32}(10^{-4})$ exactly (the exact value
$9999.99974737875$ differs from $10^4$ because $10^{-4}$ is not representable
in float, D-004), whereas the float accumulator is wrong by more than
1\,\%.

\paragraph{Dense $4\times4$ kernels.}
Block inversion uses LU with partial pivoting in double; a pivot with
magnitude below $10^{-30}$ is shifted by $10^{-30}$ and counted, so a singular
block never produces NaN or Inf. Block products that feed stored results
accumulate in double; the matrix--vector kernels inside Krylov and smoother
loops run in float.

\subsection{Krylov solvers}\label{sec:krylov}

All solvers act on the $4N$-vector with the block matrix--vector product
(native LDU traversal, processor exchange posted before and completed after
the local work) and test convergence on
$\|\vect r\|_2/\nu_{\mathrm n}$ with the outer normalisation $\nu_{\mathrm n}$
supplied by the assembler: converged if the value is below the absolute
tolerance or below $\eta_n$ times its initial value, after at least one
iteration.

\begin{description}
\item[blockBiCGStab.] Right-preconditioned BiCGStab (Saad~\cite{Saad2003},
  Alg.~7.7, arranged as the native \code{PBiCGStab}). Breakdown guards
  $|\rho|<\VSMALL\|\vect r_0\|^2$ or $|\omega|<\VSMALL$ restart from the
  current iterate with a fresh shadow residual (at most three times).
  Best-iterate return as in Section~\ref{sec:ls}.
\item[blockGMRES.] Restarted GMRES($m$) with right preconditioning
  (Alg.~9.5): Arnoldi on $\mathsf A\mathsf M^{-1}$, only the basis
  $\mathsf V$ is stored and $\vect x\mathrel{+}=\mathsf M^{-1}(\mathsf V\vect y)$
  is formed once per cycle ($m+3$ block vectors). Correct only for a fixed
  preconditioner.
\item[blockFGMRES.] Flexible GMRES($m$)~\cite{Saad1993} (Alg.~9.6): the
  preconditioned vectors $\vect z_j=\mathsf M_j^{-1}\vect v_j$ are stored, so a
  preconditioner that changes between applications (the K-cycle, an
  iterative coarsest-level solve) is handled correctly; $2m+2$ block
  vectors, $m=10$ by default and $6$ above 40\,M cells. Arnoldi uses modified
  Gram--Schmidt with double-accumulated inner products; the Hessenberg
  matrix, the Givens rotations and the least-squares solve are in double. A
  happy breakdown $h_{j+1,j}<\VSMALL\|\vect r_0\|$ ends the cycle.
\end{description}
The combination of the K-cycle with a solver other than blockFGMRES is
rejected at start-up.

\subsection{Block additive-correction AMG}\label{sec:gamg}

\paragraph{Hierarchy.} The agglomeration is OpenFOAM's own
\code{GAMGAgglomeration} (default \code{faceAreaPair}, which uses geometry
only), obtained on the mesh and cached; with \code{mergeLevels} $=2$ two
pairwise agglomeration steps are merged into one level. Restriction
$\mathsf R_l$ sums the fine residual blocks of each coarse cell, prolongation
$\mathsf P_l=\mathsf R_l^T$ is piecewise constant, and the coarse operator is
the Galerkin product $\mathsf A_{l+1}=\mathsf R_l\mathsf A_l\mathsf P_l$
evaluated by summation~\cite{Hutchinson1986,Weiss1999}: coarse diagonal
blocks collect the fine diagonal blocks of their cells and both off-diagonal
blocks of every fine face interior to the coarse cell; coarse face blocks
collect the fine face blocks with the native face-flip map; processor
interface blocks are summed with the native patch-face restriction. Sums are
in double, results are stored in float. Every level uses the same block
matrix class on the native coarse \code{lduPrimitiveMesh}.

\paragraph{Operator complexity and coarsening ratio.} With
$C_{\mathrm{op}}=\sum_l\mathrm{nnz}_l/\mathrm{nnz}_0$ (in blocks) and the
measured global ratios $r_l=N_l/N_{l+1}$, an agglomeration is accepted if
$C_{\mathrm{op}}\le1.5$ and, for W- and K-cycles, every $r_l\ge3$; otherwise
\code{mergeLevels} is increased (up to 4) and the mesh re-agglomerated, and
if the rule still fails the run stops with the measured values. The rule
follows from the work of one cycle per finest-level work unit for a constant
ratio $r$:
\begin{equation}
  W_V=\sum_{l\ge0}r^{-l}=\frac{r}{r-1},\quad
  W_F=\sum_{l\ge0}(l+1)r^{-l}=\Bigl(\frac{r}{r-1}\Bigr)^{2},\quad
  W_W=\sum_{l\ge0}\Bigl(\frac2r\Bigr)^{l}=\frac{r}{r-2}\ \ (r>2),
  \label{eq:cyclecost}
\end{equation}
see Table~\ref{tab:cyclecost}. With plain pairwise agglomeration ($r=2$) the
W-cycle, and a K-cycle that takes its second step on every level, cost
$O(L\,N)$ instead of $O(N)$.

\begin{table}[tbp]
\centering\small
\caption{Cycle work per finest-level work unit, equation~\eqref{eq:cyclecost}
(infinitely many levels). The K-cycle lies between V and W depending on how
often the second GCR step is taken.}
\label{tab:cyclecost}
\begin{tabular}{lrrrr}
\toprule
coarsening ratio $r$ & 2 & 3 & 4 & 8\\
\midrule
V & 2.00 & 1.50 & 1.33 & 1.14\\
F & 4.00 & 2.25 & 1.78 & 1.31\\
W & $\infty$ & 3.00 & 2.00 & 1.33\\
\bottomrule
\end{tabular}
\end{table}

\paragraph{Smoothers.} Block Gauss--Seidel,
$\vect x_P\leftarrow\mathsf D_P^{-1}(\vect b_P-\sum_N\mathsf A_{PN}\vect x_N)$,
with precomputed dense $4\times4$ inverses and the loop structure of the
native \code{GaussSeidelSmoother}; processor contributions are Jacobi-like
(values of the previous sweep). Block ILU(0) is implemented in the LDU-native
form, the block analogue of the native DILU preconditioner (D-015):
\begin{equation}
  \mathsf M=(\mathsf D^\ast+\mathsf L)\,\mathsf D^{\ast-1}(\mathsf D^\ast+\mathsf U),\qquad
  \mathsf D^\ast_u=\mathsf D_u-\sum_{f:\,u(f)=u}\mathsf L_f\,\mathsf D^{\ast-1}_{l(f)}\mathsf U_f,
\end{equation}
with one sweep $\vect x\leftarrow\vect x+\mathsf M^{-1}(\vect b-\mathsf A\vect x)$.
Only diagonal blocks are modified, so the factorisation stays inside the LDU
pattern; full ILU(0) would also update couplings between two neighbours of a
common cell, which are not faces and hence not stored. Processor interfaces
enter the ILU(0) sweep only through the residual.

\paragraph{Cycles.} Level $0$ is the finest, $L$ the coarsest,
$\mathcal S_l(\vect x,\vect b,\nu)$ denotes $\nu$ smoothing sweeps
($\nu_{\mathrm{pre}}=1$, $\nu_{\mathrm{post}}=2$, and $\nu=2$ for both on the
finest level). Algorithm~\ref{alg:cycle} gives the cycles; the F-cycle does
one W-style recursion per level followed by a V-cycle on the remaining
residual. The K-cycle of Notay and Vassilevski~\cite{NotayVassilevski2008},
which underlies AGMG~\cite{Notay2010}, replaces the coarse correction by one
or two steps of a generalised conjugate residual (GCR) method preconditioned
by the cycle on the next level (Algorithm~\ref{alg:kstep}); the second step
is taken only if the first reduced the coarse residual by less than the
factor $t=0.25$. The cycle thus behaves like a V-cycle where the coarse
correction is effective and like a W-cycle where it is not, per level and
per application. Additive-correction aggregation AMG has a comparatively
weak coarse-grid correction, which is why a plain V-cycle does not scale
optimally with it. Because the K-cycle is a nonlinear, variable
preconditioner, it requires the flexible outer Krylov method.

\begin{algo}{$\mathrm{CYCLE}(l,\vect b)$: approximate solution of
$\mathsf A_l\vect x=\vect b$ from $\vect x=0$.}{alg:cycle}
\textbf{if} $l=L$: \textbf{return} coarsest solve of $\mathsf A_L\vect x=\vect b$ \textit{(dense LU or BiCGStab)}\\
$\vect x\leftarrow\mathcal S_l(0,\vect b,\nu_{\mathrm{pre}})$; \quad $\vect r\leftarrow\vect b-\mathsf A_l\vect x$; \quad $\vect b_c\leftarrow\mathsf R_l\vect r$\\
\textbf{V:} \> $\vect e_c\leftarrow\mathrm{CYCLE}(l+1,\vect b_c)$\\
\textbf{F:} \> $\vect e_c\leftarrow\mathrm{CYCLE}_F(l+1,\vect b_c)$; \quad $\vect e_c\mathrel{+}=\mathrm{CYCLE}_V(l+1,\vect b_c-\mathsf A_{l+1}\vect e_c)$\\
\textbf{W:} \> $\vect e_c\leftarrow\mathrm{CYCLE}(l+1,\vect b_c)$; \quad $\vect e_c\mathrel{+}=\mathrm{CYCLE}(l+1,\vect b_c-\mathsf A_{l+1}\vect e_c)$\\
\textbf{K:} \> $\vect e_c\leftarrow\mathrm{KSTEP}(l+1,\vect b_c)$\\
$\vect x\leftarrow\vect x+\mathsf P_l\vect e_c$; \quad $\vect x\leftarrow\mathcal S_l(\vect x,\vect b,\nu_{\mathrm{post}})$; \quad \textbf{return} $\vect x$
\end{algo}

\begin{algo}{$\mathrm{KSTEP}(l,\vect b)$: Notay--Vassilevski K-step, GCR with at
most $k_{\max}=2$ steps; all inner products double-accumulated.}{alg:kstep}
$\vect r_0\leftarrow\vect b$\\
$\vect z_1\leftarrow\mathrm{CYCLE}(l,\vect r_0)$; \quad $\vect q_1\leftarrow\mathsf A_l\vect z_1$\\
$\alpha_1\leftarrow\langle\vect q_1,\vect r_0\rangle/\langle\vect q_1,\vect q_1\rangle$; \quad
$\vect e\leftarrow\alpha_1\vect z_1$; \quad $\vect r_1\leftarrow\vect r_0-\alpha_1\vect q_1$\\
\textbf{if} $\|\vect r_1\|\le t\,\|\vect r_0\|$ or $k_{\max}=1$: \textbf{return} $\vect e$\\
$\vect z_2\leftarrow\mathrm{CYCLE}(l,\vect r_1)$; \quad $\vect q_2\leftarrow\mathsf A_l\vect z_2$\\
$\theta\leftarrow\langle\vect q_2,\vect q_1\rangle/\langle\vect q_1,\vect q_1\rangle$ \quad \textit{(with the unorthogonalised $\vect q_2$)}\\
$\vect q_2\leftarrow\vect q_2-\theta\,\vect q_1$; \quad $\vect z_2\leftarrow\vect z_2-\theta\,\vect z_1$\\
$\alpha_2\leftarrow\langle\vect q_2,\vect r_1\rangle/\langle\vect q_2,\vect q_2\rangle$; \quad $\vect e\leftarrow\vect e+\alpha_2\vect z_2$; \quad \textbf{return} $\vect e$
\end{algo}

The K-step workspace ($\vect z_1,\vect q_1,\vect z_2,\vect q_2,\vect r$) is
allocated once per level and cached. On the finest level the Krylov
acceleration is provided by the outer FGMRES, so K-steps are applied on
the coarse levels only.

\paragraph{Coarsest level.} If the coarsest level resides on one rank with
at most 64 cells, it is solved by a dense LU factorisation in double of the
assembled $(4n)\times(4n)$ matrix, refactorised only when the matrix
changes. Otherwise block BiCGStab with the block-diagonal preconditioner is
used to a relative tolerance of $10^{-3}$ or 50 iterations. The selection is
automatic and logged.

\paragraph{Processor agglomeration.} On many ranks the cost of coarse levels
is latency, not arithmetic. Coarse levels are gathered onto fewer ranks with
OpenFOAM's \code{GAMGProcAgglomeration} (\code{masterCoarsest}): a level with
fewer than $5000\,n_{\mathrm{ranks}}$ cells is agglomerated onto
$\max(1,n_{\mathrm{ranks}}/4)$ ranks, and a level below 5000 cells onto one
rank. Block interface coefficients follow the native mapping with 16
coefficients per face instead of one. Skipping coarse levels is not used: it
removes the low-frequency correction and saves less than 5\,\% of the
cycle work.

\paragraph{Cycle-efficiency controller.} The preconditioner efficiency
$\rho=\|\vect b-\mathsf A\mathsf M^{-1}\vect b\|/\|\vect b\|$ is measured on the
first preconditioner application of every linear solve (in FGMRES it comes
from the first Arnoldi step at the cost of one extra reduction). Every 50
outer iterations the median $\rho_{\mathrm{med}}$ of the window is evaluated:
if $\rho_{\mathrm{med}}>0.7$ the post-smoothing count is increased (up to 4)
and, at the maximum, V is promoted to F and F to W (refused if the
coarsening-ratio rule does not hold); a K-cycle is not promoted further
(it adapts internally). If $\rho_{\mathrm{med}}<0.3$ the post-smoothing count
is decreased, and a cycle promoted by the controller is demoted by one step
once the count is back at one sweep and $\rho_{\mathrm{med}}<0.2$. Every
change requires the condition in two consecutive windows, is applied between
linear solves only, and is logged as an event; $\rho_{\mathrm{med}}>0.9$ for
two windows at the maximum setting with a W- or K-cycle is reported as a
failure.

% =============================================================================
\section{Implementation in OpenFOAM}\label{sec:impl}
% =============================================================================

\paragraph{Library structure.} coupledFoam is an add-on library
\code{libcoupledFoam} and an application built with \code{wmake} into the
user's library and application directories; OpenFOAM itself is not modified.
The library has five parts:
\code{blockMatrix} (precision types, $4\times4$ kernels, double reductions,
the block LDU matrix and processor interfaces),
\code{blockSolvers} (Krylov solvers, block-diagonal and block-GAMG
preconditioners, smoothers),
\code{assembly} (the two-stage assembler, deferred-correction convection,
Rhie--Chow interpolation, limited non-orthogonal correction, boundary
classification and implicit MRF),
\code{control} (PTC and CFL control, line search, remediation, sentinel,
convergence monitor, Anderson acceleration, GAMG controller) and
\code{io} (restart state, run information, JSON output). Project code is
compiled with \code{-Wall -Wextra -Werror} (OpenFOAM headers as system
headers) and \code{-O3 -march=native -fno-math-errno}; \code{-ffast-math} is
excluded because it breaks NaN and Inf detection.

\paragraph{Runtime selection and settings.} Solvers, preconditioners and
smoothers are selected at run time from the \code{coupled} entry of
\code{fvSolution.solvers}; the outer-iteration controls live in the
\code{coupled} sub-dictionary of \code{fvSolution}. Every constant is a
dictionary entry with a documented default; the solver prints the complete
effective settings at start-up, a single machine-parseable \code{CF|} line
per iteration (CFL, $\omega$, cuts, $R$, $r_U$, $r_p$, linear iterations and
residual, assembly/solve/turbulence times, set sizes, rollbacks, clamps) and
a JSON summary at the end (iterations, wall and CPU time summed over ranks,
peak RSS, GAMG levels and $C_{\mathrm{op}}$).

\paragraph{Parallel interfaces.} Each processor patch is wrapped by a block
interface that exchanges 4-vectors per face with non-blocking MPI, using the
native neighbour rank, tag, communicator and face ordering, on the finest
level and on the native coarse GAMG interfaces. Unlike the native
\code{interfaceBouCoeffs}, whose product is subtracted, the block
coefficients are stored as the actual matrix entries $\mathsf A_{PN}$ of the
remote neighbour. Processor-cyclic, cyclic and cyclicAMI interfaces are not
block-coupled (Section~\ref{sec:bc}).

\paragraph{Floating-point safety.} Development and test runs trap
floating-point exceptions (\code{FOAM\_SIGFPE}, \code{FOAM\_SETNAN}); any trap is
a test failure. Benchmark runs disable traps and enable flush-to-zero and
denormals-are-zero. Every division by a computed quantity, every pivot and
every argument of \code{sqrt}, \code{log} and \code{pow} is guarded; the
guards carry a common source tag so that they can be reviewed. Narrowing to
float clamps and counts, and the sentinel prevents non-finite states from
propagating or being written.

\paragraph{Implementation status.} The baseline (Sections~\ref{sec:disc}
to~\ref{sec:restart}, block BiCGStab/GMRES with the V-cycle block GAMG) was
completed first; the components of amendment B (FGMRES, Eisenstat--Walker
forcing, F/W/K cycles with the coarsening-ratio rule, dense coarsest LU,
processor agglomeration, the cycle controller, Anderson acceleration and
zonal factors) were integrated afterwards. Every run records its effective
settings, and the results in Sections~\ref{sec:vv} and~\ref{sec:perf} are
reported with the configuration that produced them.

\paragraph{Testing strategy.} Three layers, all driven by \code{pytest} and
writing one JSON record per test:
(i) unit-test applications on the real mesh classes:
\code{Test-precision} (type sizes), \code{Test-block4Ops} (1000 random
inversions, $\|\mathsf A\mathsf A^{-1}-\mathsf I\|_\infty<10^{-5}$, guarded
singular block), \code{Test-blockMatrix} (block product against native
\code{lduMatrix::Amul} componentwise, 1 and 4 ranks),
\code{Test-doubleReduce}, \code{Test-blockGAMG} (block Poisson system,
iteration count and 1-vs-4-rank identity, and the comparison of the V-, F-,
W- and K-cycles, Figure~\ref{fig:cycles}) and \code{Test-blockFGMRES}
(convergence with a deliberately variable preconditioner, where standard
GMRES may stall);
(ii) case tests T0--T5 (T5 deferred, D-063), T-restart, T-fpe and the informational T-mrf, each
against \code{simpleFoam} on the identical mesh and schemes;
(iii) the benchmark harness of Section~\ref{sec:perf}. Test runs enable
floating-point traps and record the machine load, since the development
machine is shared (D-007).

% =============================================================================
\section{Verification and validation}\label{sec:vv}
% =============================================================================

\subsection{Test cases}

Unless stated otherwise, every case is run on 1 and on 4 ranks (scotch
decomposition); both must pass, with a cross-rank tolerance of $10^{-4}$
relative on integral quantities. Native references are produced with
\code{simpleFoam} on the identical mesh, schemes and boundary conditions.

\begin{description}
\item[T0, lid-driven cavity (laminar).] $128\times128$ cells, $Re=100$ and
  $1000$, \code{linearUpwind} convection. Pass: $R<10^{-8}$ within 300
  iterations; centreline profiles $u(y)$ and $v(x)$ (129 points) within
  $10^{-3}$ relative $L_2$ of \code{simpleFoam}; no floating-point trap; no
  clamped coefficients. Ghia et al.~\cite{Ghia1982} provide the classical
  reference data.
\item[T1, pitzDaily (SST).] About 12\,k cells. Pass: $R<10^{-6}$ within 400
  iterations; inlet-to-outlet pressure drop within 1\,\% of \code{simpleFoam};
  CFL $\ge100$ without cuts after iteration 100; no rollbacks.
\item[T2, backward-facing step 2D (SST).] Pass: reattachment length (first
  sign change of the lower-wall shear stress downstream of the step) within
  2\,\% of \code{simpleFoam}; at least $2\times$ fewer outer iterations to
  $10^{-5}$ (definitions of the two residuals in D-024). Experimental
  reference: Driver and Seegmiller~\cite{Driver1985}.
\item[T3, airFoil2D (SST and GEKO).] Pass: $C_l$, $C_d$ within 0.5\,\% of
  \code{simpleFoam} for each model with identical solver settings; empty
  dynamic set at convergence.
\item[T4, motorBike.] \code{snappyHexMesh} meshes of about 350\,k cells (a)
  and 1--2\,M cells (b; retargeted from the specified 3--5\,M cells,
  D-031), 10 ranks (D-031). Each mesh is built once and reused by all tests and
  benchmarks. Pass: both solvers reach a stationary window mean of $C_d$
  and $C_l$ (averaged force criterion for oscillating wakes, D-042,
  Section~\ref{sec:perf}); window means within $\max(2\,\%,0.002)$ ($C_d$)
  and $\max(2\,\%,0.01)$ ($C_l$) of \code{simpleFoam}; RMS of the
  mean-field differences at most 0.02 (proposed, D-042 addendum); static set
  at most 1\,\% of the cells; no rollbacks.
\item[T5, Ahmed body, $25^\circ$ slant.] Generated geometry, half model,
  5--8\,M cells, $U_\infty=\SI{40}{m/s}$, SST with wall functions. Pass:
  as T4 (averaged force criterion, D-042); comparison with the experimental
  $C_d=0.285$~\cite{Ahmed1984} ($\pm10$\,\%) is informational.
  \emph{Deferred:} T5 is not run in the present test campaign (decision of
  the project owner, D-063); the case template and its mesh generation are
  kept, and no T5 results are reported in this document.
\item[T-restart.] T1 and T3-SST, split at iteration 150: iteration count
  $\pm2$, final quantities within $10^{-5}$, flux consistency below
  $10^{-6}$.
\item[T-fpe.] T1 from $\Uv=0$, no potential initialisation, $\mathrm{CFL}_0=200$,
  no upwind start-up: no trap, recovery (rollbacks allowed) and convergence
  to $R<10^{-5}$.
\item[T-mrf.] \code{mixerVessel2D} (MRF, SST), torque against
  \code{simpleFoam}; informational, since no threshold is specified (D-009).
\item[T-scaling.] Strong scaling on T4a with 1, 2, 4, 8, 12 and 16 ranks on
  physical cores, 150 iterations per run, instead of the specified T4b with
  300 iterations (12-hour budget of the final re-run, D-059; recorded as
  \code{specRanks=false}). Pass: the parallel efficiency of coupledFoam at
  the largest rank count is at least 80\,\% of that of \code{simpleFoam}
  (unchanged threshold; Figure~\ref{fig:scaling}).
\end{description}

\subsection{Results}

Table~\ref{tab:tests} lists all test records produced so far.

\cftable{tests}{Test results.}{tab:tests}

\paragraph{T0, $Re=100$.} coupledFoam converged to $R<10^{-8}$ in
\cfnum{TZeroReOneZeroZeroIterations} outer iterations; the reference
\code{simpleFoam} (SIMPLEC) run needed \cfnum{TZeroReOneZeroZeroNativeIterations}
iterations to its residual tolerance of $10^{-8}$ (the two residual
normalisations differ, Section~\ref{sec:norms}). The centreline profiles agree
to a relative $L_2$ difference of \cfsci{TZeroReOneZeroZeroLTwou} in $u$ and
\cfsci{TZeroReOneZeroZeroLTwov} in $v$ (Figure~\ref{fig:t0profiles}), far
below the pass threshold of $10^{-3}$ and consistent with two converged
solutions of the same discretisation. The history
(Figure~\ref{fig:t0history}) shows the CFL ramp and the line-search factor.
The corresponding profiles at $Re=1000$ are shown in
Figure~\ref{fig:t0profilesRe1000}; the test record is listed in
Table~\ref{tab:tests}.
The wall times of the test records were measured on a shared machine and are
not benchmark results.

\cffigure{T0_Re100_profiles}{T0, lid-driven cavity at $Re=100$: centreline
profiles $u(y)$ and $v(x)$, coupledFoam against simpleFoam.}{fig:t0profiles}

\cffigure{T0_Re100_np1_history}{T0, $Re=100$, one rank: combined residual $R$
of coupledFoam and native initial residuals of simpleFoam (left); CFL number
and line-search factor $\omega$ (right).}{fig:t0history}

\cffigure{T0_Re1000_profiles}{T0, lid-driven cavity at $Re=1000$: centreline
profiles, coupledFoam against simpleFoam.}{fig:t0profilesRe1000}

\paragraph{T1--T4.} Integral quantities and their differences to the native
solver are collected in Table~\ref{tab:validation}; remediation activity on
the snappyHexMesh cases is shown in Figure~\ref{fig:remediation} and
Table~\ref{tab:remcat}. T5 is deferred (D-063).

\cftable{validation}{Validation, T1--T4: integral quantities (pressure drop,
reattachment length, $C_l$, $C_d$, torque) of coupledFoam and simpleFoam and
their relative differences.}{tab:validation}

\paragraph{Force histories.} Figures~\ref{fig:ovWing}
and~\ref{fig:ovMotorbike} overlay the drag and lift coefficients of both
solvers over the wall-clock time for the wing (T3) and the motorbike (T4a,
T4b), with the running mean over the averaging window of the D-042 rule and the
final window mean; the complete per-iteration histories of every run,
including residuals, cost per iteration and the linear-solver controls, are
in Appendix~\ref{app:histories}. Histories of runs that are still being
computed end early and are completed when the report is regenerated.

\cffigure{loads_overview_wing}{Wing (T3 airFoil2D, SST and GEKO): $C_d$
(left) and $C_l$ (right) over wall-clock time, simpleFoam (orange) and
coupledFoam (blue); dashed: running mean over the D-042 window, dotted:
final window mean; light band: per-iteration minimum-maximum per group of
iterations; a second row per case zooms on the shorter run where the wall
times differ by more than 5x.}{fig:ovWing}
\cffigure{loads_overview_motorbike}{motorBike (T4a, T4b): $C_d$ and $C_l$
over wall-clock time, both solvers; light band: per-iteration
minimum-maximum per group of iterations; a second row per case zooms on the
shorter run where the wall times differ by more than 5x.}{fig:ovMotorbike}

\cffigure{remediation_history}{Size of the dynamic remediation set over the
iterations, one panel per run.}{fig:remediation}

\cftable{remediation_categories}{Remediation cells per category and test
run.}{tab:remcat}


\clearpage
\subsection{Flow fields}\label{sec:flowfields}

Integral quantities can agree through a compensation of errors, so the
converged fields of the two solvers are also compared directly. For every
2D case with a coupledFoam run and a \code{simpleFoam} reference on the same
mesh, the figures of this section show the velocity magnitude
$|\Uv|/U_{ref}$ with streamlines and the pressure coefficient
$C_p=(p-p_0)/(\tfrac12U_{ref}^2)$ of both solvers on one colour scale, and
the pointwise difference coupledFoam$-$\code{simpleFoam} on a diverging
colour scale. The fields are read from the latest time directory of each
run with the OpenFOAM reader of VTK, interpolated to the mesh points and cut
at mid-depth (\code{bench/plot\_fields2d.py}); the meshes are identical, so
the difference is evaluated node by node, and the quoted RMS is weighted
with the area. $U_{ref}$ is the lid speed (T0), the inlet speed (T1:
\SI{10}{m/s}, T2: \SI{44.2}{m/s}) or the free-stream speed (T3:
\SI{26}{m/s}); $p_0$ is the outlet pressure, except in the closed cavity,
where the area-weighted mean pressure is removed from both fields.

\paragraph{T0.} The velocity differences are at the level of the converged
profile differences: RMS \cfnum{FldTZeroReOneZeroZeroURms} at $Re=100$ and
\cfnum{FldTZeroReOneZeroZeroZeroURms} at $Re=1000$
(Figures~\ref{fig:fldTzeroU}--\ref{fig:fldTzerothousandP}). The largest
differences sit in the two upper corners, where the moving lid meets the
walls: the pressure is singular there and both solutions are least
resolved.

\cffigure{fields_T0_Re100_U}{T0, $Re=100$: $|\Uv|/U_{lid}$ with streamlines,
coupledFoam (left) and simpleFoam (middle) on one colour scale, and the
difference (right).}{fig:fldTzeroU}
\cffigure{fields_T0_Re100_p}{T0, $Re=100$: pressure coefficient (mean
removed) of both solvers and the difference.}{fig:fldTzeroP}
\cffigure{fields_T0_Re1000_U}{T0, $Re=1000$: $|\Uv|/U_{lid}$ with
streamlines and the difference.}{fig:fldTzerothousandU}
\cffigure{fields_T0_Re1000_p}{T0, $Re=1000$: pressure coefficient and the
difference.}{fig:fldTzerothousandP}

\clearpage
\paragraph{T1.} The recirculation behind the step, its length and the
velocity profiles coincide (Figures~\ref{fig:fldToneU}--\ref{fig:prTone});
the RMS differences are \cfnum{FldTOneURms} in $|\Uv|/U_{ref}$ and
\cfnum{FldTOnePRms} in $C_p$, consistent with the pressure-drop difference
of Table~\ref{tab:validation}.

\cffigure{fields_T1_U}{T1 pitzDaily: $|\Uv|/U_{ref}$ with streamlines,
coupledFoam (top), simpleFoam (middle), difference (bottom).}{fig:fldToneU}
\cffigure{fields_T1_p}{T1 pitzDaily: pressure coefficient and the
difference.}{fig:fldToneP}
\cffigure{profiles_T1}{T1 pitzDaily: streamwise velocity profiles at the
stations $x/h$ downstream of the step ($h$ step height), simpleFoam thick
orange, coupledFoam dashed blue.}{fig:prTone}

\clearpage
\paragraph{T2.} Separation at the step corner, the recirculation bubble and
the recovery agree (Figures~\ref{fig:fldTtwoU}--\ref{fig:prTtwo}); the RMS
difference in $|\Uv|/U_{ref}$ is \cfnum{FldTTwoURms}, the maximum
(\cfnum{FldTTwoUMax}) sits in the thin shear layer leaving the step corner.
The lower-wall skin friction crosses zero at the same point, which is the
reattachment-length test. The \code{simpleFoam} skin friction has a narrow
spike near $x/h\approx1.4$ that coupledFoam does not have; the
\code{simpleFoam} run did not reach the residual target in 20\,000
iterations (Table~\ref{tab:tests}), so the spike is not interpreted
further.

\cffigure{fields_T2_U}{T2 backward-facing step: $|\Uv|/U_{ref}$ with
streamlines, coupledFoam (top), simpleFoam (middle), difference
(bottom).}{fig:fldTtwoU}
\cffigure{fields_T2_p}{T2 backward-facing step: pressure coefficient and
the difference.}{fig:fldTtwoP}
\cffigure{profiles_T2}{T2 backward-facing step: velocity profiles at the
stations $x/h$ (top) and lower-wall skin friction $C_f$ (bottom); the zero
crossing downstream of the step is the reattachment point.}{fig:prTtwo}

\clearpage
\paragraph{T3.} With GEKO the fields agree away from the airfoil (RMS
\cfnum{FldTThreeGEKOURms} in $|\Uv|/U_{ref}$), but the suction peak of
coupledFoam is higher and the suction-side pressure plateau near
$x/c\approx0.1$--$0.2$ is missing (Figure~\ref{fig:prTthreegeko}); this is
the source of the lift difference of Table~\ref{tab:validation}. With SST
the coupledFoam run of the test record did not converge (3000 iterations,
Table~\ref{tab:tests}); its field shows a large separated region with
vortices above the suction side that the \code{simpleFoam} solution does not
have (Figure~\ref{fig:fldTthreesstU}; RMS difference
\cfnum{FldTThreeKOmegaSSTURms}). The figures show the state of the test
records at the time of writing, not a converged comparison.

\cffigure{fields_T3_GEKO_U}{T3 airFoil2D, GEKO: $|\Uv|/U_\infty$ with
streamlines around the airfoil, coupledFoam (top), simpleFoam (middle),
difference (bottom).}{fig:fldTthreegekoU}
\cffigure{fields_T3_GEKO_p}{T3 airFoil2D, GEKO: pressure coefficient and
the difference.}{fig:fldTthreegekoP}
\cffigure{profiles_T3_GEKO}{T3 airFoil2D, GEKO: surface pressure
coefficient $-C_p$ over the chord.}{fig:prTthreegeko}
\cffigure{fields_T3_kOmegaSST_U}{T3 airFoil2D, SST: $|\Uv|/U_\infty$ with
streamlines; the coupledFoam run of this record did not
converge.}{fig:fldTthreesstU}
\cffigure{fields_T3_kOmegaSST_p}{T3 airFoil2D, SST: pressure coefficient
and the difference.}{fig:fldTthreesstP}
\cffigure{profiles_T3_kOmegaSST}{T3 airFoil2D, SST: surface pressure
coefficient.}{fig:prTthreesst}

\clearpage
\paragraph{T4 (3D).} The motorbike cases are rendered with ParaView
(\code{bench/render\_fields.py}, \code{pvbatch}, decomposed cases read in
place): the body surface mesh, $|\Uv|/U_\infty$ and $C_p$ on the vertical
mid-plane $y=y_c$ and on a horizontal plane at 20\,\% of the body height,
and $C_p$ on the body, each for \code{simpleFoam} and coupledFoam with
fixed colour ranges ($|\Uv|/U_\infty\in[0,1.3]$, $C_p\in[-1.5,1]$,
identical for both solvers). Panels whose run has not finished are drawn
as ``results pending''. The Ahmed body (T5) is not part of the present test
campaign by decision of the project owner (D-063); no T5 results are
reported.

\paragraph{Mean-field differences.} Side-by-side absolute fields hide
differences of a few per cent of $U_\infty$, and the instantaneous wake does
not settle (D-042). The decisive field comparison is therefore between the
window-mean fields $\overline{\Uv}$, $\overline{p}$ of both solvers
(coupledFoam: the last $W$ iterations of the run; the cached
\code{simpleFoam} reference: its continuation, Section~\ref{sec:perf}),
and it is shown as a signed
difference, coupledFoam minus \code{simpleFoam}, computed cell by cell on
the identical decomposed mesh by \code{coupledFieldCompare}:
\begin{equation}
  \frac{\Delta\overline{u}_x}{U_\infty}
    =\frac{\overline{u}_{x,\mathrm{cf}}-\overline{u}_{x,\mathrm{sf}}}{U_\infty},
  \qquad
  \Delta\overline{C}_p
    =\frac{\overline{p}_{\mathrm{cf}}-\overline{p}_{\mathrm{sf}}}{\tfrac12U_\infty^2},
  \label{eq:meandelta}
\end{equation}
on the two planes and, with the boundary values, on the body surface. Both
use a diverging colour map with fixed symmetric limits $\pm0.2$, white at
zero, so that equal colours mean equal differences in every figure and in
every run. Red marks regions where coupledFoam has the higher mean
streamwise velocity or pressure, blue where it has the lower. A wake that is
shorter or longer in one solver appears as a pair of opposite-signed bands
behind the body; a shifted separation line on the body appears as a
red--blue dipole of $\Delta\overline{C}_p$ across that line. Differences
confined to the shear layers of the wake are expected, because the window
means of two oscillating solutions converge slowly there; differences on
the body surface of the size of the $C_d$ difference in
Table~\ref{tab:validation} identify where that difference originates. The
volume RMS of both quantities is the field criterion of the validation
table (proposed limit 0.02).

\cffigure{render_T4a_geometry}{T4a motorBike, 354\,k cells: body surface
mesh.}{fig:rTfouraGeo}
\cffigure{render_T4a_slice_U}{T4a: $|\Uv|/U_\infty$ on the mid-plane (top)
and the horizontal plane (bottom), simpleFoam (left) and coupledFoam
(right).\cfrenderflag{TFoura}}{fig:rTfouraU}
\cffigure{render_T4a_slice_p}{T4a: $C_p$ on the mid-plane and the horizontal
plane.\cfrenderflag{TFoura}}{fig:rTfouraP}
\cffigure{render_T4a_surface_Cp}{T4a: surface $C_p$ on the
body.\cfrenderflag{TFoura}}{fig:rTfouraCp}
\cffigure{render_T4a_mean_U}{T4a: window-mean velocity magnitude on the
mid-plane, both solvers.\cfrenderflag{TFoura}}{fig:rTfouraMean}
\cffigure{render_T4a_delta_slices}{T4a: mean-field differences
coupledFoam$-$\code{simpleFoam}, Eq.~\eqref{eq:meandelta}: streamwise
velocity $\Delta\overline{u}_x/U_\infty$ (top) and $\Delta\overline{C}_p$
(bottom) on the mid-plane (left) and the horizontal plane (right); limits
$\pm0.2$, white: no difference.\cfrenderflag{TFoura}}{fig:rTfouraDelta}
\cffigure{render_T4a_delta_surface}{T4a: mean surface-pressure difference
$\Delta\overline{C}_p$ (coupledFoam$-$\code{simpleFoam}) on the body,
oblique and top view; limits $\pm0.2$.\cfrenderflag{TFoura}}{fig:rTfouraDeltaS}
\clearpage
\IfFileExists{figures/render_T4b_delta_slices.png}{%
\cffigure{render_T4b_geometry}{T4b motorBike, 1.70\,M cells: body surface
mesh.}{fig:rTfourbGeo}
\cffigure{render_T4b_slice_U}{T4b: $|\Uv|/U_\infty$ on the mid-plane (top)
and the horizontal plane (bottom), simpleFoam (left) and coupledFoam
(right).\cfrenderflag{TFourb}}{fig:rTfourbU}
\cffigure{render_T4b_slice_p}{T4b: $C_p$ on the mid-plane and the horizontal
plane.\cfrenderflag{TFourb}}{fig:rTfourbP}
\cffigure{render_T4b_surface_Cp}{T4b: surface $C_p$ on the
body.\cfrenderflag{TFourb}}{fig:rTfourbCp}
\cffigure{render_T4b_mean_U}{T4b: window-mean velocity magnitude on the
mid-plane, both solvers.\cfrenderflag{TFourb}}{fig:rTfourbMean}
\cffigure{render_T4b_delta_slices}{T4b: mean-field differences
coupledFoam$-$\code{simpleFoam}, Eq.~\eqref{eq:meandelta}:
$\Delta\overline{u}_x/U_\infty$ (top) and $\Delta\overline{C}_p$ (bottom) on
both planes; limits $\pm0.2$.\cfrenderflag{TFourb}}{fig:rTfourbDelta}
\cffigure{render_T4b_delta_surface}{T4b: mean surface-pressure difference
$\Delta\overline{C}_p$ on the body, oblique and top view; limits
$\pm0.2$.\cfrenderflag{TFourb}}{fig:rTfourbDeltaS}
\clearpage}{%
\paragraph{T4b.} No mean-field comparison of the fine motorbike mesh is
available for this document: T4b is run only after T4a has been accepted
(D-063), and its figures are added when the report is regenerated after
that run.}
\clearpage

% =============================================================================
\section{Performance}\label{sec:perf}
% =============================================================================

\subsection{Benchmark protocol}

The goal is the solver that reaches a converged answer in the least time,
not the one with the fewest iterations. The primary metrics are therefore the
wall-clock time and the CPU time (CPU-hours summed over all ranks) to
convergence, measured for all solvers with the same criterion.

\paragraph{Configurations.} For each case in \{T1, T2, T3-SST, T3-GEKO, T4a,
T4b\}, with identical mesh, decomposition, schemes, turbulence model and
boundary conditions, the configurations of Table~\ref{tab:configs} are run
(D-025, amendment B10). Configuration~A is the tutorial as shipped with
OpenFOAM: pitzDaily, backwardFacingStep2D and motorBike use SIMPLEC
(\code{consistent yes}) with their own relaxation factors; the airFoil2D
tutorial uses plain SIMPLE with $p$ relaxed by 0.3 and $\Uv$ and the
turbulence equations by 0.7, so A on T3 overrides the case setting
\code{consistent yes} (which serves the reference run of
Section~\ref{sec:vv}) with \code{consistent no}. The Ahmed body (T5) has no
tutorial; A there means the motorBike tutorial settings, the closest
external-aerodynamics tutorial. As in the motorBike tutorial script,
configuration~A on the motorbike runs \code{potentialFoam} before
\code{simpleFoam}, and its cost is counted as pre-processing of~A; on all
other cases, and for B everywhere, the native configurations start from the
initial fields of the case. Configurations D--H isolate single components of
coupledFoam; their case scope is given in Table~\ref{tab:configs}.
Since D-043 the coupledFoam defaults (C) are a V-cycle with the damped block
ILU(0) smoother and the cycle-efficiency controller off. The configuration
of amendment B10 that fixed the V-cycle (E: \code{cycleType V},
\code{autoTune no}) is therefore identical to C and is no longer run;
configuration H sets \code{cycleType K} explicitly (with the controller
off) and is the cycle comparison, and H-tune adds the controller to the
K-cycle, which isolates the controller (D-068). The E-* variants of
amendment C7 keep their names and change one setting of C each. The
harness records the final cycle type and post-smoothing count of every
coupledFoam run, so that a comparison is only reported between runs that
actually differ.

\paragraph{Scope of the present campaign (D-063).} Not every case receives
every configuration. The light cases T1--T3 carry most of the benchmark;
T4a runs the full matrix; T4b is run with few configurations and only after
T4a has been accepted; the single-precision variants
(Section~\ref{sec:sp}) run on a few cases only. The Ahmed body T5 is
deferred by decision of the project owner: the configurations listed for
T5 in Table~\ref{tab:configs} and in the specification are not run, and
no T5 entry appears in the results. Strong scaling is measured on T4a
(D-059).

\begin{table}[tbp]
\centering\small
\caption{Benchmark configurations (D-025, amendment B10).}
\label{tab:configs}
\begin{tabular}{llp{6.2cm}p{2.6cm}}
\toprule
 & solver & settings & cases (D-063)\\
\midrule
A & \code{simpleFoam} & tutorial settings (see text; T3: SIMPLE, $p$ 0.3,
  $\Uv$, $k$, $\omega$ 0.7; motorbike: \code{potentialFoam} start) &
  T1--T3, T4a\\
B & \code{simpleFoam} & SIMPLEC, relaxation $p$ 1.0, $\Uv$ 0.9, $k,\omega$ 0.9
  (tuned native baseline; identical to A on T1) & T2, T3, T4a\\
C & coupledFoam & defaults (V-cycle, damped block ILU(0), controller off,
  Eisenstat--Walker; D-043) & T1--T3, T4a, T4b\\
D & coupledFoam & block-diagonal preconditioner (isolates the AMG) & T1--T3\\
F & coupledFoam & fixed relative tolerance, \code{adaptiveRelTol no} &
  T1, T3-SST\\
G & coupledFoam & Anderson acceleration, \code{anderson.enabled yes} &
  T1, T3-SST\\
H & coupledFoam & fixed K-cycle, \code{cycleType K}, \code{autoTune no} &
  T1--T3, T4a\\
H-tune & coupledFoam & K-cycle with the controller, \code{cycleType K},
  \code{autoTune yes} (the default before D-043) & T1, T2\\
E-* & coupledFoam & C with one change each (amendment C7): scalar
  Rhie--Chow coefficient, static non-orthogonality threshold 60 or
  $65^\circ$, algebraic pair agglomeration, SFD off, $\eta_{\max}=0.7$ &
  T2 or T3 (one case each)\\
\bottomrule
\end{tabular}
\end{table}

\paragraph{Stop criterion.} The solvers' own stop criteria are disabled and
a fixed iteration budget is run. The harness determines the convergence
iteration identically for all solvers: over the last 100 iterations,
$\max C_d-\min C_d\le0.002\,|\overline{C_d}|$ and the same for $C_l$ (the
pressure drop for T1 and T2).

\paragraph{Wake cases (D-042 and addendum).} T4a, T4b and T5 have physically
oscillating wakes: the T4b \code{simpleFoam} reference oscillates by about
$\pm1\,\%$ in $C_d$ and $\pm8\,\%$ in $C_l$ indefinitely, so the window
criterion above cannot be met by any steady solver on these cases. For them
(user decisions, D-042 and its addenda) a run is converged when its window
mean is stationary: over a window of the last $W$ iterations the means of
the two half-windows differ by at most $\max(0.01\,|\overline{C}|,0.005)$,
for $C_d$ and $C_l$ separately. The convergence iteration is the first
iteration (scanned in steps of 50) at which this holds, and the coefficients
compared are the window means, within $\max(2\,\%,0.002)$ for $C_d$ and
$\max(2\,\%,0.01)$ for $C_l$. The window is \emph{common to both solvers of
a case} (D-068; $W=400$ on T4a and T4b). The earlier rule derived the window
from each run's own iteration budget $n$ ($W=\max(300,n/2)$, D-042 addenda);
because the scan cannot start before iteration $W$, the longer native
budgets alone then delayed the earliest possible convergence of
\code{simpleFoam} (e.g.\ $W=1500$ for a 3000-iteration reference against
$W=400$ for an 800-iteration coupledFoam run), and the resulting time to
convergence was set by the budgets rather than by the flow. The per-run
window is kept as a sensitivity value only. With either window the window
means of the references agree to within a few tenths of a per cent. Note
that the absolute floor 0.005 of the stationarity tolerance is 6--8\,\% of
$C_l$ on the motorbike, so that the stationarity test is weak for $C_l$.
Numbers alone can agree for the wrong reason, so the flow fields are
compared as well. A \code{fieldAverage} function object averages $\Uv$ and
$p$; for coupledFoam it covers the last $W$ iterations of the run. The
cached \code{simpleFoam} references were run before the averaging existed
and were continued once, with the averaging active from the first continued
iteration; their mean fields therefore cover the continuation (1500
iterations on T4a, 2000 on T4b), a longer window at a later stage of the
same stationary oscillation. The utility \code{coupledFieldCompare} writes
the differences of the mean fields between the two solvers, including the
boundary values (the surface $\Delta C_p$ on the body), and reports their
volume-weighted RMS. The proposed pass thresholds are an RMS of
$|\Delta\overline{\Uv}|/U_\infty$ and of
$|\Delta\overline{p}|/(\tfrac12U_\infty^2)$ of at most 0.02 each. The
criteria of T0--T3 are unchanged.

\paragraph{Timing.} Every rank of every application that the case script
starts runs under \code{/usr/bin/time -v}. The total cost of a solution is
the wall time (maximum over ranks) and the CPU time (user plus system,
summed over ranks, in CPU-hours) of the solver \emph{plus} that of its
pre-processing: the \code{potentialFoam} initialisation, which the case
script runs for every coupledFoam run and for configuration~A on the
motorbike (tutorial start). The
pre-processing runs before the solver, so its wall and CPU times are added;
both parts are also recorded separately in every result. The time to
convergence is
\begin{equation}
  t_{\mathrm{conv}} = t_{\mathrm{pre}} + f\,t_{\mathrm{solver}},\qquad
  c_{\mathrm{conv}} = c_{\mathrm{pre}} + f\,c_{\mathrm{solver}},
\end{equation}
with $t$ the wall and $c$ the CPU time, and $f$ the solver's own wall-clock
progress at the convergence iteration (coupledFoam's cumulative wall time per
iteration, \code{simpleFoam}'s \code{ClockTime}) divided by its total
(D-025). The pre-processing is a fixed offset and is not scaled. Using the
wall-clock fraction $f$ also for the CPU time assumes a constant CPU-to-wall
ratio during the run. The time per iteration (Figure~\ref{fig:breakdown}) is
that of the solver alone. Peak memory is the per-rank maximum RSS, reported
as sum and maximum over ranks. On T1--T3 each configuration is run three
times and the median (with minimum and maximum) is reported; the motorbike
configurations are run once (12-hour budget, D-059). Speed-up is reported both as
a wall-time and as a CPU-hour ratio against the better of A and B. Results
are tagged with a hash of their configuration; results of an outdated
configuration are excluded from all tables and figures.

\paragraph{Acceptance of the B10 variants.} For F, the adaptive tolerance of
C must not be more than 5\,\% slower than the fixed tolerance, in wall time
and in CPU-hours, and the monitored quantity ($C_d$, or $\Delta p$ for T1)
at the end of the fixed iteration budget must agree to $10^{-4}$ (relative).
G and H have no threshold; their differences to C are reported in wall
time and CPU-hours.

\paragraph{Hardware.} 16 physical cores (32 logical with SMT, one socket),
94\,GB RAM visible to WSL. The test and benchmark runs of the motorbike use
10 ranks (D-031), the strong-scaling test up to 16 (D-059); all timing runs
use physical cores only, with MPI ranks bound to cores. The benchmarks are
run when no other solver runs on the machine, and the load average before
each run is recorded (D-007).

\subsection{Results}

The executive summary (Table~\ref{tab:summary}) compares the best native
configuration with coupledFoam in wall-clock time and CPU-hours; the
per-configuration results are in Figures~\ref{fig:benchwall} (wall time)
and~\ref{fig:benchcpu} (CPU-hours), the iteration counts in
Figure~\ref{fig:benchiters}. Figure~\ref{fig:breakdown} splits the time per
coupledFoam iteration, and Figure~\ref{fig:memory} compares the measured
peak memory with the budget of Table~\ref{tab:memory}. The acceptance of the
B10 variants is evaluated in Table~\ref{tab:btenacc}, with the cycle
comparison (H and H-tune against C) in Figure~\ref{fig:benchcycles} and Anderson
acceleration (G) in Figure~\ref{fig:anderson}. The block-GAMG hierarchies
are listed in Table~\ref{tab:gamglevels}, the Eisenstat--Walker forcing terms
and preconditioner efficiencies are shown in Figure~\ref{fig:etarho}, and the
strong scaling in Figure~\ref{fig:scaling}.

\paragraph{Convergence point chosen by the user.} The automatic criteria
decide whether a test passes. For the force cases (T3, T4) the user may in
addition name an earlier iteration from which they consider a run converged,
after looking at the load histories of Appendix~\ref{app:histories}
(\texttt{report/user\_convergence.json}, D-060). Times, CPU-hours, speed-ups and
coefficients are then taken at that iteration; Table~\ref{tab:convchoice} lists
both points. In the load histories the automatic point is dash-dotted, the user
point solid.

\cftable{convergence_choice}{Convergence point of the force cases: automatic
and user choice (D-060).}{tab:convchoice}

\cftable{executive_summary}{Executive summary: time to convergence (identical
window criterion, coupledFoam including its potentialFoam initialisation),
best native configuration against coupledFoam, as wall-clock time and
CPU-hours.}{tab:summary}

The fine motorbike T4b is benchmarked with configuration C only (D-059,
D-063); its comparison with \code{simpleFoam} is the test run against the
cached reference in Table~\ref{tab:wakespeed}
(Section~\ref{sec:speedtests}), together with the conditions of that
reference.

\cffigure{bench_wall}{Wall-clock time to convergence (median of the repeats:
three on T1--T3, one on the motorbike) per case and
configuration.}{fig:benchwall}

\cffigure{bench_cpu}{CPU time to convergence (CPU-hours summed over ranks,
median of the repeats).}{fig:benchcpu}

\cffigure{bench_iters}{Outer iterations to convergence. Shown for
completeness; the decisive metrics are those of Figures~\ref{fig:benchwall}
and~\ref{fig:benchcpu}.}{fig:benchiters}

\cffigure{bench_breakdown}{coupledFoam time per outer iteration split into
assembly, linear solve and turbulence: wall time (left) and CPU time summed
over ranks (right, split in the proportions of the wall-time
breakdown).}{fig:breakdown}

\cffigure{bench_memory}{Peak memory, sum of the per-rank maximum RSS, with
the B7 budget of Table~\ref{tab:memory} (105--120\,GB at 45\,M cells) scaled
to the cell count of each case (grey band). On small meshes the fixed
per-process overhead of OpenFOAM dominates, so the comparison is meaningful
only for the larger cases.}{fig:memory}

\cftable{b10_acceptance}{Amendment B10: configurations F, G and H against
C in wall time and CPU-hours; pass/fail for F.}{tab:btenacc}

\cffigure{bench_cycles}{Benchmark cycle comparison: default V-cycle (C),
fixed K-cycle (H) and K-cycle with autoTune (H-tune), wall time and
CPU-hours to convergence.}{fig:benchcycles}

\cffigure{anderson}{Anderson acceleration on and off (T1, T3-SST): outer
iterations, wall time and CPU-hours to convergence.}{fig:anderson}

\cffigure{scaling}{Strong scaling on T4a with 1, 2, 4, 8, 12 and 16 ranks (D-059):
wall time per iteration, parallel efficiency and CPU-hours of the fixed
iteration count, coupledFoam and simpleFoam. Pass: the parallel efficiency
of coupledFoam at the largest rank count is at least 80\,\% of that of
simpleFoam.}{fig:scaling}

\cftable{gamg_levels}{Block-GAMG hierarchy per run: cycle type and
post-smoothing sweeps at start and end (autoTune), levels, operator
complexity, cells, ranks and coarsening ratios per level, coarsest-level
solver.}{tab:gamglevels}

\cffigure{cycle_comparison}{Unit test \code{Test-blockGAMG} (not a
benchmark): Krylov iterations, solve wall time and process CPU time of the
V-, F-, W- and K-cycles on the block-Poisson system of the $128\times128$
cavity mesh, serial, solved to $10^{-9}$ (D-023).}{fig:cycles}

\cffigure{eta_rho_history}{Eisenstat--Walker forcing term $\eta_n$ and
preconditioner efficiency $\rho$ with the controller events, one row per
run.}{fig:etarho}

\clearpage
\subsection{Speed-up measured on the test runs}\label{sec:speedtests}

Until the benchmark matrix has run, the test records give a first
measurement of the speed-up (single runs on the shared machine, not
benchmark quality; D-007): serial on T0--T3, 10 ranks on the motorbike.
Both solvers are evaluated at the same criterion: the residual target of
the test ($10^{-8}$ for T0, $10^{-5}$ for T1 and T2), reached by the
combined residual $R$ of coupledFoam and by all initial residuals of
\code{simpleFoam} (on T2 the native iteration count of D-024); the
stationary window mean with the common window on the motorbike (D-068);
on T3 the force window of criterion 12.3(ii) on both force histories.
Wall-clock time and CPU-hours are taken at that iteration from each
solver's time log (D-068, harness review M1). The residual normalisations differ (Section~\ref{sec:norms}), so the time
axis, not the residual level, is the comparable quantity; the field
comparisons of Section~\ref{sec:flowfields} show that both runs end at the
same solution. On the serial runs CPU-hours and wall-clock time are nearly
proportional; both are reported.

Figure~\ref{fig:speedres} plots the residuals against wall-clock time.
Figure~\ref{fig:speedconv} gives the time and CPU-hours to convergence:
coupledFoam is faster by a factor of \cfnum{SpeedWallTZeroReOneZeroZero}
in wall-clock time (\cfnum{SpeedCpuTZeroReOneZeroZero} in CPU-hours) on T0
at $Re=100$, \cfnum{SpeedWallTZeroReOneZeroZeroZero}
(\cfnum{SpeedCpuTZeroReOneZeroZeroZero}) at $Re=1000$ and
\cfnum{SpeedWallTOne} (\cfnum{SpeedCpuTOne}) on T1
(\cfnum{SpeedWallCfTOne}\,s against \cfnum{SpeedWallSfTOne}\,s). On T2
\code{simpleFoam} did not reach the target within
\cfnum{SpeedItersSfTTwo} iterations, so the factor
\cfnum{SpeedWallTTwo} (\cfnum{SpeedCpuTTwo} in CPU-hours) is a lower
bound. \cfnum{SpeedNotConvNote}
Figure~\ref{fig:speediter} shows why the iteration count alone is
misleading: one coupled iteration costs \cfnum{SpeedIterCostRatioTTwo}
times a \code{simpleFoam} iteration on T2 and
\cfnum{SpeedIterCostRatioTThreeGEKO} times on T3-GEKO, and the linear
solve takes \cfnum{SpeedSolveShareTTwo}\,\% (T2) to
\cfnum{SpeedSolveShareTThreeGEKO}\,\% (T3-GEKO) of the iteration.
Figure~\ref{fig:speedtrade} places both effects in one picture.

\cffigure{speed_residual_wall}{Residual against wall-clock time (symmetric
log time axis): coupledFoam combined residual $R$ and simpleFoam initial
residuals of $p$ and $U_x$. Normalisations differ; the time axis is the
comparable quantity. Cases whose coupledFoam run is not finished show the
reference only.}{fig:speedres}

\cffigure{speed_time_to_conv}{Wall-clock time (left) and CPU-hours (right)
to convergence, simpleFoam (orange) and coupledFoam (blue), both timed to
the same criterion: T0--T2 the test's residual target (coupledFoam $R$;
simpleFoam every initial residual below the same value), T3 the 12.3(ii)
force window, T4 the first stationary D-042 window mean with one
400-iteration window for both solvers (D-068). T0--T3 serial, solver time
only; T4 on 10 MPI ranks, coupledFoam including potentialFoam, against the
cached simpleFoam reference (run without the tutorial's potentialFoam
start, timing conditions not recorded). Numbers: speed-up
simpleFoam/coupledFoam. Hatched: criterion not reached, bar = whole run
(simpleFoam: lower bound, marked $\geq$).}{fig:speedconv}

\paragraph{Motorbike: coupledFoam against the cached references.} On the
motorbike the benchmark runs few configurations (D-059, D-063), and the
speed-up is taken from the test records instead: the coupledFoam test run
against the cached \code{simpleFoam} reference of the same mesh and
decomposition, both to the stationary window mean with the common window
(D-068). Table~\ref{tab:wakespeed} gives the result together with the
conditions that limit its fairness: the reference is a single native
configuration (the tutorial relaxation, not the better of A and B), it may
have been started without the tutorial's \code{potentialFoam}
initialisation, and the machine load during the cached reference runs was
not recorded. The speed-up with each run's own window (the earlier rule) is
given as a sensitivity value only; it is not a result, because that rule
made the earliest possible convergence point of the longer reference run
depend on its iteration budget. The mean fields of the references, used by
the field comparison, cover their continuation of 1500 (T4a) or 2000 (T4b)
iterations. Result: \cfnum{WakeSpeedNote}

\cftable{wake_speedup}{Motorbike speed-up from the test records against
the cached simpleFoam references.}{tab:wakespeed}

\cffigure{speed_iteration_cost}{Time per outer iteration. Left: share of
assembly, linear solve and turbulence in a coupledFoam iteration. Right:
wall time per iteration of both solvers and the cost ratio
coupledFoam/simpleFoam.}{fig:speediter}

\cffigure{speed_tradeoff}{Iterations against wall-clock time. Each arrow
leads from simpleFoam to coupledFoam for one case; diagonals are lines of
constant time per iteration. coupledFoam trades far fewer iterations for
more expensive ones; it is faster where the arrow points
down.}{fig:speedtrade}

\clearpage
\subsection{Metrics for comparison with other solvers}\label{sec:metrics}

Validation and performance reports of commercial and research coupled
solvers~\cite{cmp-Fluent12TG,cmp-StarEffortless2020,cmp-HelyxCoupled,Uroic2021}
typically quantify cost per unit of mesh, memory per million cells,
residual reduction against iterations and time, parallel efficiency, and
the growth of the iteration count with the mesh size. The same quantities
are computed here from the existing test records and logs, so that the
results of this paper can be placed next to such reports; they are no
substitute for a comparison on identical cases and hardware.
Table~\ref{tab:metcost} normalises the cost per outer iteration and the peak
memory by the cell count, Table~\ref{tab:metres} gives the iterations and
wall-clock time to reduce the residual by two to four orders of magnitude,
Table~\ref{tab:metscal} the parallel efficiency of the strong-scaling test,
and Table~\ref{tab:metgrid} the dependence on the mesh of the motorbike.
Robustness on a deliberately degraded mesh, weak scaling, energy to
solution and a memory breakdown by data structure are not reported: they
require a degraded-mesh variant of a test case, a mesh series at constant
cells per rank, a power measurement, and allocation counters in the solver,
none of which is part of the present campaign.

\cftable{metrics_cost}{Cost and memory per million cells.}{tab:metcost}
\cftable{metrics_residual}{Residual reduction: iterations and wall-clock
time.}{tab:metres}
\cftable{metrics_scaling}{Strong scaling: time per iteration and parallel
efficiency.}{tab:metscal}
\cftable{metrics_grid}{Mesh dependence on the motorbike (T4a against
T4b).}{tab:metgrid}

\clearpage
\subsection{Single precision}\label{sec:sp}

Amendment D adds full single precision (SP) as a build option: a second
OpenFOAM v2606 build with \code{WM\_PRECISION\_OPTION=SP}, in which fields,
geometry, all matrices and all solves are 4-byte floats. The reference is
the double-precision build (DP) used for every other result of this paper;
the amendment names the mixed build \code{SPDP} as the reference, but no
such build is made (D-062). Both builds use identical compiler flags (the
unmodified \code{linux64Gcc} rules, \code{-O3}), so that the timings compare
precision and not compilers. The meshes are generated in DP and copied to
the SP case in ASCII with 12 significant digits; SP is a solve-only mode.

\paragraph{Where the gain comes from.} coupledFoam already stores its block
matrix and all Krylov and multigrid vectors in float
(Section~\ref{sec:storage}). SP halves in addition the fields, the mesh
geometry, the scalar matrices of the turbulence equations and the files
written, so its expected gain is confined to assembly, turbulence and I/O,
about 15--20\,\% of the wall-clock time according to the amendment (D12).
\code{simpleFoam}, whose linear solvers work in double in the DP build, is
run in SP as well. In the benchmark the SP configurations are F1
(\code{simpleFoam}, settings of B) and F2 (coupledFoam, settings of C);
each solver is compared with itself in the other precision, never across
precisions.

\paragraph{Safeguards.} Everything accumulated over many cells stays in
double in both builds: norms, dot products, sums, the Galerkin coarse
operators, the small dense systems of the Krylov and Anderson methods, and
the force and moment integration (a dedicated double-accumulated force
evaluation, used identically for both solvers). The stabilising constants
\code{SMALL} and \code{VSMALL}, whose values change with the precision, are
replaced by explicit guard constants. Two geometric risks are specific to SP.
A float resolves about seven significant digits, so a coordinate of 10\,m
is represented to about 1\,$\mu$m; the difference of two nearby cell
centres, from which face weights, gradients and wall distances are formed,
loses the digits that the absolute coordinates carry. Before every SP run
the mesh is therefore translated so that the centre of its bounding box is
the origin; the centre of rotation of the moments, probes and sampling
lines are shifted with it, which is where an omission would produce silently
wrong moments. The shifted mesh must then pass a checkMesh gate: no check
may fail in SP that passes in DP, and no cell may have a negative volume or
a wrong orientation; otherwise the case is marked \code{SP-geometry-fail}
and its SP benchmark is skipped. The static remediation cells are
recomputed in SP, and the difference in their number is reported. Finally,
the precision profile raises the attainable targets in SP (outer residual
target $10^{-5}$ instead of $10^{-6}$, absolute linear tolerance $10^{-6}$,
$\eta_{\min}=10^{-2}$), because a field increment below about $10^{-6}$ of
the field value is lost to rounding; the force window remains the decisive
convergence test.

\paragraph{Evaluation.} SP is compared with DP on a few cases only
(D-063). Table~\ref{tab:spperf} gives the wall-clock time, CPU-hours and
peak memory of both solvers in both precisions,
Table~\ref{tab:spbreak} the time per coupledFoam iteration by component,
Figure~\ref{fig:spfloor} the residual histories (the convergence floor),
and Table~\ref{tab:spverdict} the verdict per case: SP is \emph{usable} on a
case if the monitored quantity of coupledFoam ($C_d$; $\Delta p$ for T1,
$x_r/h$ for T2) lies within 0.5\,\% of DP and the checkMesh gate passed,
otherwise \emph{not usable}, with the reason. \cfnum{SpStatus}

\cftable{sp_performance}{Single against double precision: wall-clock time,
CPU-hours and peak memory of both solvers.}{tab:spperf}
\cftable{sp_breakdown}{coupledFoam time per iteration by component, DP and
SP.}{tab:spbreak}
\cffigure{sp_convergence_floor}{Convergence floor: combined residual $R_n$
of coupledFoam in DP (solid) and SP (dashed) on the same case; dotted: the
SP residual target of the precision profile.}{fig:spfloor}
\cftable{sp_verdict}{Single-precision verdict per case: monitored quantity
in DP and SP, checkMesh gate, change of the static remediation set, origin
shift, verdict.}{tab:spverdict}

% =============================================================================
\section{Discussion and limitations}\label{sec:discussion}
% =============================================================================

\subsection{What single precision requires}

Four observations made during development determine whether a float block
solver reaches double-precision steady residuals.
(i) The residual must be exact: $b-Ax$ is formed in double from the double
assembly (D-012). With that, the float matrix only affects the convergence
rate, as in mixed-precision iterative refinement.
(ii) Near-singular systems must be made nonsingular exactly. A weak pressure
reference leaves a near-null space in which the float solve drifts once the
right-hand side is rounded; the strong reference of D-021 removes it without
approximation.
(iii) Linear solves near the float floor can fail. A best-iterate return
(D-020) and treating a solve that misses its target as a failed step with a
CFL cut (D-029) keep a converged state from being damaged.
(iv) Tolerances must be consistent with the float floor and with the outer
target. The absolute linear tolerance must lie below the outer target, and at
least one Krylov iteration must always be made (D-022). Conversely, the
float floor limits how closely two solves can agree: two solves of the block
Poisson test stopped at $\|\vect r\|_2/\nu_{\mathrm n}=10^{-8}$ differ by
\cfsci{ExpDTwoThreeRelDiffTolEightMinus} (relative $L_2$) between 1 and 4
ranks, \cfsci{ExpDTwoThreeRelDiffTolNineMinus} at $10^{-9}$, and tighter
tolerances do not reduce the difference below about
\cfsci{ExpDTwoThreeRelDiffFloor} (D-023; values from the table of that
entry, the logs were not kept). Cross-rank identity of float solves must therefore
be tested at a tolerance matched to that floor.

The precision types are the library's own, because OpenFOAM's \code{SPDP}
build is the inverse of what a memory-bound block solver needs (D-001).

\subsection{Deferred correction at high CFL}\label{sec:dcstab}

With frozen flux, exact linear solves and $\omega=1$, one outer iteration
with deferred correction is the stationary iteration
$\vect x_{k+1}=\vect x_k+\mathsf M^{-1}(\vect b-\mathsf A_{\mathrm{HO}}\vect x_k)$ with
$\mathsf M=\mathsf A_{\mathrm{UD}}+\mathsf T$, $\mathsf T=\operatorname{diag}(V/\dt)$.
Its error propagates with
\begin{equation}
  \mathsf G = \mathsf I-(\mathsf A_{\mathrm{UD}}+\mathsf T)^{-1}\mathsf A_{\mathrm{HO}}
  = (\mathsf A_{\mathrm{UD}}+\mathsf T)^{-1}\bigl(\mathsf T-\mathsf C\bigr),
  \qquad \mathsf C=\mathsf A_{\mathrm{HO}}-\mathsf A_{\mathrm{UD}},
  \label{eq:dcG}
\end{equation}
and converges only if $\rho(\mathsf G)<1$. The pseudo-time term plays two
roles: it damps $\mathsf C$ relative to $\mathsf T$, and it slows the
smooth modes. As CFL$\to\infty$, $\mathsf T\to0$ and
$\mathsf G\to-\mathsf A_{\mathrm{UD}}^{-1}\mathsf C$, so the damping by the
pseudo-time term disappears exactly when the CFL control removes it.

For a first estimate, the symbol of~\eqref{eq:dcG} was evaluated for constant
convection at angle $\theta$ on a uniform Cartesian grid, with OpenFOAM's
\code{linearUpwind} (face value from the upwind cell plus its central
gradient, which in 1D is Fromm's scheme), first-order upwind in
$\mathsf A_{\mathrm{UD}}$, central diffusion at the cell P\'eclet number
$Pe_h$, and $\mathsf T$ from~\eqref{eq:dt} with $V^{1/3}=h$; modes with
wavelength longer than the domain ($|\vect k|<2\pi/(128\,h)$) are excluded,
all other wave vectors in $[-\pi/h,\pi/h]^2$ are sampled. The script
\code{bench/exploratory\_numbers.py} computes the values quoted here. For 1D
convection the symbol of
$\mathsf A_{\mathrm{UD}}^{-1}\mathsf C$ is $\tfrac{\mathrm i}{2}\sin\xi$, so
$\rho\le\tfrac12$ even at infinite CFL. In 2D, cross-stream modes whose
upwind operator consists almost entirely of numerical diffusion, which the
correction removes, bring $\rho$ close to one. In the limit CFL$\to\infty$ at
$\theta=45^\circ$ the result is $\rho(\mathsf G)=\cfnum{ExpSymRhoInfPeLow}$ for
$Pe_h=\cfnum{ExpSymPeLow}$ (the $Re=1000$ cavity on $128^2$ cells, based on
the lid velocity) and $\rho(\mathsf G)=\cfnum{ExpSymRhoInfPeHigh}$ for
$Pe_h=\cfnum{ExpSymPeHigh}$. At CFL~500 and $Pe_h=\cfnum{ExpSymPeLow}$ the
value rises to \cfnum{ExpSymRhoCflFiveHundredPeLow}: at finite CFL the
largest $|\mathsf G|$ comes from the smoothest resolvable modes, which the
pseudo-time term slows, not from the correction. In this idealised model, therefore, $\rho<1$ at every CFL,
but with a small margin at high cell P\'eclet numbers. Effects outside the
model can push the iteration over the edge: non-uniform and distorted cells,
boundaries and the singular corners of the cavity, the Picard lag of the
convecting flux, and above all an inexact linear solve, which replaces
$\mathsf M^{-1}$ by a different approximate inverse.

This is consistent with the observations on the $Re=1000$ cavity in
exploratory runs (logs in \code{run/exp\_re1000}, outside the test protocol).
With the solver default of that time (BiCGStab preconditioned by block GAMG
with a Gauss--Seidel smoother, relative tolerance
\cfnum{ExpReOneZeroZeroZeroDefaultRelTol}), the run stagnated once deferred
correction was switched on after the upwind start-up: after the start-up the
CFL number stayed between \cfnum{ExpReOneZeroZeroZeroDefaultCflMinAfterStartup}
and \cfnum{ExpReOneZeroZeroZeroDefaultCflMaxAfterStartup} (5th to 95th
percentile \cfnum{ExpReOneZeroZeroZeroDefaultCflPZeroFiveAfterStartup} to
\cfnum{ExpReOneZeroZeroZeroDefaultCflPNineFiveAfterStartup}) with
\cfnum{ExpReOneZeroZeroZeroDefaultCflCutsTotal} CFL cuts, and after
\cfnum{ExpReOneZeroZeroZeroDefaultIterations} iterations $R$ was
\cfsci{ExpReOneZeroZeroZeroDefaultFinalR} (minimum
\cfsci{ExpReOneZeroZeroZeroDefaultMinR}). With pure upwind throughout, the
same configuration reached CFL~500 at iteration
\cfnum{ExpReOneZeroZeroZeroUpwindFirstIterAtCFLmax}
(\cfnum{ExpReOneZeroZeroZeroUpwindCflCutsTotal} CFL cuts in total) and
$R<10^{-8}$ in \cfnum{ExpReOneZeroZeroZeroUpwindFirstIterBelowTarget}
iterations. With the block ILU(0) smoother and otherwise identical settings,
deferred correction reached CFL~500 at iteration
\cfnum{ExpReOneZeroZeroZeroIluZeroFirstIterAtCFLmax} with
\cfnum{ExpReOneZeroZeroZeroIluZeroCflCutsTotal} cuts and converged to
$R<10^{-8}$ in \cfnum{ExpReOneZeroZeroZeroIluZeroFirstIterBelowTarget}
iterations. The stagnation is thus not
intrinsic to deferred correction at CFL~500 for this case; it arises from a
small theoretical margin combined with an inaccurate inner solve. The design
consequences are the Eisenstat--Walker forcing (tighter solves as the outer
iteration converges), the stronger smoothers and cycles of
Section~\ref{sec:gamg}, the treatment of failed solves, and remediation,
which sets $\beta_P=0$ where the correction is locally destabilising. A
rigorous statement for the actual discretisation would require the spectrum
of~\eqref{eq:dcG} for the assembled operators, which has not been computed.

On the $Re=100$ cavity an exploratory comparison of linear solvers (logs in
\code{run/exp\_linsolver}) showed the practical weight of the inner solve:
BiCGStab with the Gauss--Seidel smoother needed
\cfnum{ExpLinBicgGsIterations} outer iterations and
\cfnum{ExpLinBicgGsLinIters} Krylov iterations in total, with
\cfnum{ExpLinBicgGsMaxed} solves reaching the iteration limit, whereas the
ILU(0) smoother needed \cfnum{ExpLinBicgIluIterations} outer and
\cfnum{ExpLinBicgIluLinIters} Krylov iterations and reduced the summed solve
time by a factor of \cfnum{ExpLinSolveRatioGsIlu}
(\cfnum{ExpLinBicgGsTSolve}\,s against \cfnum{ExpLinBicgIluTSolve}\,s);
GMRES with the Gauss--Seidel smoother was similarly fast
(\cfnum{ExpLinGmresGsTSolve}\,s, \cfnum{ExpLinGmresGsIterations} outer
iterations). These runs were made on a
loaded machine and are indicative only; the benchmark of
Section~\ref{sec:perf} is authoritative.

\subsection{Other limitations}

\begin{itemize}
\item \emph{Pseudo-time in the Rhie--Chow coefficient.} $D=V/\bar a$ includes
  the pseudo-time term, so the converged solution depends weakly on the
  final local time step through the Rhie--Chow dissipation, as SIMPLE
  solutions depend on the relaxation factor. At the large CFL numbers
  reached at convergence the effect is small; the cavity comparison
  (Section~\ref{sec:vv}) bounds it for that case.
\item \emph{Explicit cyclic and AMI coupling.} Cyclic, cyclicAMI and
  processor-cyclic interfaces are lagged. Periodic cases converge more
  slowly; implicit block coupling across them is Phase~2.
\item \emph{Per-cell gradient limiting} applies to the Rhie--Chow pressure
  gradient only; the gradient inside the user's convection scheme cannot be
  limited per cell (D-018).
\item \emph{Boundary conditions.} Outflow through a fixed-value velocity face
  uses the native boundary value (D-019); \code{inletOutlet} and
  \code{totalPressure} are linearised about the previous outer iteration.
\item \emph{MRF} is implicit but gated by an informational test only
  (D-009).
\item \emph{Atomic write} is implemented at file level (the restart
  dictionary is renamed last) and skipped with the collated file handler
  (D-017).
\item \emph{Scope and hardware.} The 45\,M-cell production case and its
  memory budget cannot be verified on the development machine (94\,GB, D-007);
  the largest mesh of the present campaign is T4b (T5 is deferred, D-063).
\item \emph{Residual definitions} of coupledFoam and \code{simpleFoam} differ;
  the iteration criterion of T2 is defined in D-024, and all timing
  comparisons use the solver-independent force window (for the wake cases
  T4 and T5 the stationary window mean, D-042).
\end{itemize}

\subsection{Phase 2}

Block-selective AMG~\cite{Uroic2018,Uroic2021} as an additional preconditioner;
adaptive mesh refinement with rebalancing (the restart format already
reserves the refinement history); a single-precision PETSc field-split
reference; implicit cyclic and AMI block coupling; GPU execution; and the
45\,M-cell production run with its memory budget verified by peak RSS.

% =============================================================================
\section{Conclusions}\label{sec:concl}
% =============================================================================

coupledFoam solves the steady incompressible RANS equations in OpenFOAM
with a fully coupled $4\times4$-block pressure--velocity system, implicit
boundary conditions through the native patch linearisation, implicit MRF,
and a continuity row that is by construction the mass balance of the output
flux. Storing the block system and all Krylov and multigrid vectors in single
precision while forming the residual and all reductions in double turns the
outer iteration into a mixed-precision defect correction: on the cavity test
the solver converges to $R<10^{-8}$ and reproduces the native solution to
\cfsci{TZeroReOneZeroZeroLTwou}. Robustness comes from pseudo-transient
continuation with local, solution-limited time steps, a physicality line
search, cell-set remediation, rollback and the explicit handling of failed
linear solves; efficiency is sought from a block additive-correction AMG with
K-cycles under flexible GMRES and Eisenstat--Walker forcing. The exploratory
cavity runs show that the accuracy of the inner solve decides whether
deferred correction survives at high CFL numbers, which is why the linear
solver and the globalisation are designed together. Whether the coupled
solver is faster than tuned native SIMPLEC in wall-clock time and CPU-hours
is answered by the benchmark of Section~\ref{sec:perf}, whose results are
generated directly from the measurement records.

\paragraph{Reproducibility.} All measured numbers, tables and figures are
generated by \code{bench/make\_report.py} from the JSON records in
\code{results/}; the numbers of the exploratory runs and of the symbol
analysis in Section~\ref{sec:discussion} are extracted or computed by
\code{bench/exploratory\_numbers.py} (with the source of every value
recorded in \code{results/exploratory/}). The title page states the commit
(\code{\cfnum{Commit}}) at which they were generated.

% =============================================================================
\clearpage
\appendix
\section{Iteration histories}\label{app:histories}
% =============================================================================

This appendix shows, for every test run, everything that changes over the
outer iterations, so that convergence can be judged by eye and the cost of
each iteration can be traced (generated by \code{bench/plot\_histories.py}
from the solver logs, the \code{forceCoeffs} and \code{surfaceFieldValue}
outputs and the run summaries). Each run has four figures: the monitored
quantities (loads) of both solvers with the running mean and $\pm$RMS band
over the averaging window $W$ of the D-042 rule (common to both solvers of a
case, D-068), the shaded
final window and the first iteration with a stationary window (dash-dot);
the residuals of both solvers on the same y range; the wall-clock time per
iteration, cumulative with a linear fit, and the coupledFoam split into
assembly, linear solve, turbulence and other; and the linear solver and
pseudo-time controls. The left columns use a logarithmic iteration axis when
the two runs differ in length by more than a factor of three. Quantities
that stay constant over a run are listed in the text instead of plotted.
Runs that are still being computed are skipped and appear as ``results
pending''; they are filled in when the report is regenerated. The simpleFoam
reference of the 4-rank runs is the serial reference run.

% \cfhistfigure (used by the generated figures/histories.tex, which only
% \providecommand's it): [H] placement, but the first figure after the
% bullet list is shrunk to the space left on the page if at least 40 % of
% the text height is left, so that no page holds only the bullets
\makeatletter
\newlength\cf@hh
\newcommand{\cfhistfigure}[3]{%
  \par
  \ifdim\pagegoal=\maxdimen
    \setlength\cf@hh{0.8\textheight}%
  \else
    % reserve the typeset caption plus the float and caption skips
    \setbox\@tempboxa\vbox{\hsize=\linewidth\normalsize Figure 999: #2}%
    \setlength\cf@hh{\dimexpr\pagegoal-\pagetotal-\ht\@tempboxa
      -\dp\@tempboxa-4\baselineskip\relax}%
    \ifdim\cf@hh>0.8\textheight\setlength\cf@hh{0.8\textheight}\fi
    \ifdim\cf@hh<0.4\textheight\setlength\cf@hh{0.8\textheight}\fi
  \fi
  \begin{figure}[H]
    \centering
    \IfFileExists{figures/#1.pdf}%
      {\includegraphics[width=0.95\linewidth,height=\cf@hh,keepaspectratio]{figures/#1.pdf}}%
      {\pending{Figure \texttt{figures/\detokenize{#1}} is written by
        \texttt{bench/make\_report.py} once its data exist.}}
    \caption{#2}\label{#3}
  \end{figure}}
\makeatother
\IfFileExists{figures/histories.tex}{% Generated by bench/plot_histories.py - do not edit
\providecommand{\cfhistfigure}[3]{%
  \begin{figure}[H]
    \centering
    \IfFileExists{figures/#1.pdf}%
      {\includegraphics[width=0.95\linewidth,height=0.8\textheight,keepaspectratio]{figures/#1.pdf}}%
      {\pending{Figure \texttt{figures/\detokenize{#1}} is written by
        \texttt{bench/make\_report.py} once its data exist.}}
    \caption{#2}\label{#3}
  \end{figure}}
\subsection{T0 cavity, Re 100, 1 rank}\label{app:hist:T0_Re100_np1}

\begin{itemize}\setlength{\itemsep}{0pt}\small
\item coupledFoam: 57 iterations, 5.742\,s wall-clock, 0.00171 CPU-h; time per iteration 0.0567\,s (first 10\,\%) and 0.0591\,s (last 10\,\%), linear fit 0.104\,s/it ($R^2=0.976$).
\item simpleFoam: 1932 iterations, 173\,s wall-clock, 0.0481 CPU-h; time per iteration 0.0597\,s (first 10\,\%) and 0.0697\,s (last 10\,\%), linear fit 0.0913\,s/it ($R^2=0.999$).
\item coupledFoam linear iterations in total: 465.
\item constant over the whole coupledFoam run (panel omitted): bounding events = 0, cuts = 0, nRollback = 0.
\item no monitored integral quantity in this case; convergence is judged by the residuals and the final centreline profiles.
\end{itemize}

\cfhistfigure{hist_T0_Re100_np1_residuals}{T0 cavity, Re 100, 1 rank: residuals per outer iteration (left) and over wall-clock time (right), coupledFoam top, simpleFoam bottom, same y range.}{fig:hist:T0_Re100_np1:residuals}
\cfhistfigure{hist_T0_Re100_np1_wall}{T0 cavity, Re 100, 1 rank: cumulative wall time with linear fit (top), wall time per iteration with rolling median (middle), coupledFoam time split into assembly, linear solve, turbulence and other (bottom).}{fig:hist:T0_Re100_np1:wall}
\cfhistfigure{hist_T0_Re100_np1_linear}{T0 cavity, Re 100, 1 rank: linear solver and pseudo-time controls per outer iteration; bottom: simpleFoam linear iterations per field.}{fig:hist:T0_Re100_np1:linear}
\clearpage

\subsection{T0 cavity, Re 100, 4 ranks}\label{app:hist:T0_Re100_np4}

\begin{itemize}\setlength{\itemsep}{0pt}\small
\item coupledFoam: 57 iterations, 5.723\,s wall-clock, 0.00692 CPU-h; time per iteration 0.072\,s (first 10\,\%) and 0.112\,s (last 10\,\%), linear fit 0.105\,s/it ($R^2=0.967$).
\item simpleFoam: 1932 iterations, 173\,s wall-clock, 0.0481 CPU-h; time per iteration 0.0597\,s (first 10\,\%) and 0.0697\,s (last 10\,\%), linear fit 0.0913\,s/it ($R^2=0.999$).
\item coupledFoam linear iterations in total: 382.
\item constant over the whole coupledFoam run (panel omitted): bounding events = 0, cuts = 0, nRollback = 0.
\item no monitored integral quantity in this case; convergence is judged by the residuals and the final centreline profiles.
\end{itemize}

\cfhistfigure{hist_T0_Re100_np4_residuals}{T0 cavity, Re 100, 4 ranks: residuals per outer iteration (left) and over wall-clock time (right), coupledFoam top, simpleFoam bottom, same y range.}{fig:hist:T0_Re100_np4:residuals}
\cfhistfigure{hist_T0_Re100_np4_wall}{T0 cavity, Re 100, 4 ranks: cumulative wall time with linear fit (top), wall time per iteration with rolling median (middle), coupledFoam time split into assembly, linear solve, turbulence and other (bottom).}{fig:hist:T0_Re100_np4:wall}
\cfhistfigure{hist_T0_Re100_np4_linear}{T0 cavity, Re 100, 4 ranks: linear solver and pseudo-time controls per outer iteration; bottom: simpleFoam linear iterations per field.}{fig:hist:T0_Re100_np4:linear}
\clearpage

\subsection{T0 cavity, Re 1000, 1 rank}\label{app:hist:T0_Re1000_np1}

\begin{itemize}\setlength{\itemsep}{0pt}\small
\item coupledFoam: 111 iterations, 26.01\,s wall-clock, 0.00715 CPU-h; time per iteration 0.0834\,s (first 10\,\%) and 0.15\,s (last 10\,\%), linear fit 0.274\,s/it ($R^2=0.978$).
\item simpleFoam: 1971 iterations, 179\,s wall-clock, 0.0497 CPU-h; time per iteration 0.0897\,s (first 10\,\%) and 0.0797\,s (last 10\,\%), linear fit 0.0898\,s/it ($R^2=0.999$).
\item coupledFoam linear iterations in total: 1347.
\item constant over the whole coupledFoam run (panel omitted): bounding events = 0, omega = 1.
\item no monitored integral quantity in this case; convergence is judged by the residuals and the final centreline profiles.
\end{itemize}

\cfhistfigure{hist_T0_Re1000_np1_residuals}{T0 cavity, Re 1000, 1 rank: residuals per outer iteration (left) and over wall-clock time (right), coupledFoam top, simpleFoam bottom, same y range.}{fig:hist:T0_Re1000_np1:residuals}
\cfhistfigure{hist_T0_Re1000_np1_wall}{T0 cavity, Re 1000, 1 rank: cumulative wall time with linear fit (top), wall time per iteration with rolling median (middle), coupledFoam time split into assembly, linear solve, turbulence and other (bottom).}{fig:hist:T0_Re1000_np1:wall}
\cfhistfigure{hist_T0_Re1000_np1_linear}{T0 cavity, Re 1000, 1 rank: linear solver and pseudo-time controls per outer iteration; bottom: simpleFoam linear iterations per field.}{fig:hist:T0_Re1000_np1:linear}
\clearpage

\subsection{T0 cavity, Re 1000, 4 ranks}\label{app:hist:T0_Re1000_np4}

\begin{itemize}\setlength{\itemsep}{0pt}\small
\item coupledFoam: 257 iterations, 70.84\,s wall-clock, 0.079 CPU-h; time per iteration 0.0964\,s (first 10\,\%) and 0.132\,s (last 10\,\%), linear fit 0.271\,s/it ($R^2=0.961$).
\item simpleFoam: 1971 iterations, 179\,s wall-clock, 0.0497 CPU-h; time per iteration 0.0897\,s (first 10\,\%) and 0.0797\,s (last 10\,\%), linear fit 0.0898\,s/it ($R^2=0.999$).
\item coupledFoam linear iterations in total: 2763.
\item constant over the whole coupledFoam run (panel omitted): bounding events = 0, omega = 1.
\item no monitored integral quantity in this case; convergence is judged by the residuals and the final centreline profiles.
\end{itemize}

\cfhistfigure{hist_T0_Re1000_np4_residuals}{T0 cavity, Re 1000, 4 ranks: residuals per outer iteration (left) and over wall-clock time (right), coupledFoam top, simpleFoam bottom, same y range.}{fig:hist:T0_Re1000_np4:residuals}
\cfhistfigure{hist_T0_Re1000_np4_wall}{T0 cavity, Re 1000, 4 ranks: cumulative wall time with linear fit (top), wall time per iteration with rolling median (middle), coupledFoam time split into assembly, linear solve, turbulence and other (bottom).}{fig:hist:T0_Re1000_np4:wall}
\cfhistfigure{hist_T0_Re1000_np4_linear}{T0 cavity, Re 1000, 4 ranks: linear solver and pseudo-time controls per outer iteration; bottom: simpleFoam linear iterations per field.}{fig:hist:T0_Re1000_np4:linear}
\clearpage

\subsection{T1 pitzDaily, 1 rank}\label{app:hist:T1_np1}

\begin{itemize}\setlength{\itemsep}{0pt}\small
\item coupledFoam: 400 iterations, 28.1\,s wall-clock, 0.00779 CPU-h; time per iteration 0.0836\,s (first 10\,\%) and 0.0629\,s (last 10\,\%), linear fit 0.0674\,s/it ($R^2=0.992$).
\item simpleFoam: 769 iterations, 64\,s wall-clock, 0.0179 CPU-h; time per iteration 0.0698\,s (first 10\,\%) and 0.0798\,s (last 10\,\%), linear fit 0.0846\,s/it ($R^2=1.000$).
\item coupledFoam linear iterations in total: 985.
\item simpleFoam window mean of $\Delta p$: -5.7267 $\pm$ 0.011\,m$^2$/s$^2$.
\item coupledFoam window mean of $\Delta p$: -5.729 $\pm$ 0.00091\,m$^2$/s$^2$.
\item constant over the whole coupledFoam run (panel omitted): bounding events = 0, cuts = 0, nRollback = 0, omega = 1.
\end{itemize}

\cfhistfigure{hist_T1_np1_loads}{T1 pitzDaily, 1 rank: monitored quantities per outer iteration (left) and over wall-clock time (right); thin: raw, dashed: running mean over the D-042 window $W$ with $\pm$RMS band, shaded: final averaging window, dash-dot: iterations to a stationary window.}{fig:hist:T1_np1:loads}
\cfhistfigure{hist_T1_np1_residuals}{T1 pitzDaily, 1 rank: residuals per outer iteration (left) and over wall-clock time (right), coupledFoam top, simpleFoam bottom, same y range.}{fig:hist:T1_np1:residuals}
\cfhistfigure{hist_T1_np1_wall}{T1 pitzDaily, 1 rank: cumulative wall time with linear fit (top), wall time per iteration with rolling median (middle), coupledFoam time split into assembly, linear solve, turbulence and other (bottom).}{fig:hist:T1_np1:wall}
\cfhistfigure{hist_T1_np1_linear}{T1 pitzDaily, 1 rank: linear solver and pseudo-time controls per outer iteration; bottom: simpleFoam linear iterations per field.}{fig:hist:T1_np1:linear}
\clearpage

\subsection{T1 pitzDaily, 4 ranks}\label{app:hist:T1_np4}

\begin{itemize}\setlength{\itemsep}{0pt}\small
\item coupledFoam: 400 iterations, 17.53\,s wall-clock, 0.0188 CPU-h; time per iteration 0.0694\,s (first 10\,\%) and 0.0266\,s (last 10\,\%), linear fit 0.0428\,s/it ($R^2=0.958$).
\item simpleFoam: 769 iterations, 64\,s wall-clock, 0.0179 CPU-h; time per iteration 0.0698\,s (first 10\,\%) and 0.0798\,s (last 10\,\%), linear fit 0.0846\,s/it ($R^2=1.000$).
\item coupledFoam linear iterations in total: 866.
\item simpleFoam window mean of $\Delta p$: -5.7267 $\pm$ 0.011\,m$^2$/s$^2$.
\item coupledFoam window mean of $\Delta p$: -5.7296 $\pm$ 0.0022\,m$^2$/s$^2$.
\item constant over the whole coupledFoam run (panel omitted): bounding events = 0, cuts = 0, nRollback = 0, omega = 1.
\end{itemize}

\cfhistfigure{hist_T1_np4_loads}{T1 pitzDaily, 4 ranks: monitored quantities per outer iteration (left) and over wall-clock time (right); thin: raw, dashed: running mean over the D-042 window $W$ with $\pm$RMS band, shaded: final averaging window, dash-dot: iterations to a stationary window.}{fig:hist:T1_np4:loads}
\cfhistfigure{hist_T1_np4_residuals}{T1 pitzDaily, 4 ranks: residuals per outer iteration (left) and over wall-clock time (right), coupledFoam top, simpleFoam bottom, same y range.}{fig:hist:T1_np4:residuals}
\cfhistfigure{hist_T1_np4_wall}{T1 pitzDaily, 4 ranks: cumulative wall time with linear fit (top), wall time per iteration with rolling median (middle), coupledFoam time split into assembly, linear solve, turbulence and other (bottom).}{fig:hist:T1_np4:wall}
\cfhistfigure{hist_T1_np4_linear}{T1 pitzDaily, 4 ranks: linear solver and pseudo-time controls per outer iteration; bottom: simpleFoam linear iterations per field.}{fig:hist:T1_np4:linear}
\clearpage

\subsection{T2 backward-facing step, 1 rank}\label{app:hist:T2_np1}

\begin{itemize}\setlength{\itemsep}{0pt}\small
\item coupledFoam: 1000 iterations, 227.9\,s wall-clock, 0.0632 CPU-h; time per iteration 0.199\,s (first 10\,\%) and 0.194\,s (last 10\,\%), linear fit 0.224\,s/it ($R^2=0.995$).
\item simpleFoam: 20000 iterations, 1383\,s wall-clock, 0.384 CPU-h; time per iteration 0.07\,s (first 10\,\%) and 0.07\,s (last 10\,\%), linear fit 0.0693\,s/it ($R^2=1.000$).
\item coupledFoam linear iterations in total: 11796.
\item simpleFoam window mean of $\Delta p$: -39.734 $\pm$ 0.43\,m$^2$/s$^2$.
\item coupledFoam window mean of $\Delta p$: -40.188 $\pm$ 0.0064\,m$^2$/s$^2$.
\item constant over the whole coupledFoam run (panel omitted): omega = 1.
\item the reattachment length is evaluated only at the end of the run (no history); the pressure difference is shown instead.
\end{itemize}

\cfhistfigure{hist_T2_np1_loads}{T2 backward-facing step, 1 rank: monitored quantities per outer iteration (left) and over wall-clock time (right); thin: raw, dashed: running mean over the D-042 window $W$ with $\pm$RMS band, shaded: final averaging window, dash-dot: iterations to a stationary window.}{fig:hist:T2_np1:loads}
\cfhistfigure{hist_T2_np1_residuals}{T2 backward-facing step, 1 rank: residuals per outer iteration (left) and over wall-clock time (right), coupledFoam top, simpleFoam bottom, same y range.}{fig:hist:T2_np1:residuals}
\cfhistfigure{hist_T2_np1_wall}{T2 backward-facing step, 1 rank: cumulative wall time with linear fit (top), wall time per iteration with rolling median (middle), coupledFoam time split into assembly, linear solve, turbulence and other (bottom).}{fig:hist:T2_np1:wall}
\cfhistfigure{hist_T2_np1_linear}{T2 backward-facing step, 1 rank: linear solver and pseudo-time controls per outer iteration; bottom: simpleFoam linear iterations per field.}{fig:hist:T2_np1:linear}
\clearpage

\subsection{T2 backward-facing step, 4 ranks}\label{app:hist:T2_np4}

\begin{itemize}\setlength{\itemsep}{0pt}\small
\item coupledFoam: 1000 iterations, 90.4\,s wall-clock, 0.1 CPU-h; time per iteration 0.141\,s (first 10\,\%) and 0.0404\,s (last 10\,\%), linear fit 0.084\,s/it ($R^2=0.961$).
\item simpleFoam: 20000 iterations, 1383\,s wall-clock, 0.384 CPU-h; time per iteration 0.07\,s (first 10\,\%) and 0.07\,s (last 10\,\%), linear fit 0.0693\,s/it ($R^2=1.000$).
\item coupledFoam linear iterations in total: 6298.
\item simpleFoam window mean of $\Delta p$: -39.734 $\pm$ 0.43\,m$^2$/s$^2$.
\item coupledFoam window mean of $\Delta p$: -40.188 $\pm$ 0.0035\,m$^2$/s$^2$.
\item constant over the whole coupledFoam run (panel omitted): cuts = 0, nRollback = 0.
\item the reattachment length is evaluated only at the end of the run (no history); the pressure difference is shown instead.
\end{itemize}

\cfhistfigure{hist_T2_np4_loads}{T2 backward-facing step, 4 ranks: monitored quantities per outer iteration (left) and over wall-clock time (right); thin: raw, dashed: running mean over the D-042 window $W$ with $\pm$RMS band, shaded: final averaging window, dash-dot: iterations to a stationary window.}{fig:hist:T2_np4:loads}
\cfhistfigure{hist_T2_np4_residuals}{T2 backward-facing step, 4 ranks: residuals per outer iteration (left) and over wall-clock time (right), coupledFoam top, simpleFoam bottom, same y range.}{fig:hist:T2_np4:residuals}
\cfhistfigure{hist_T2_np4_wall}{T2 backward-facing step, 4 ranks: cumulative wall time with linear fit (top), wall time per iteration with rolling median (middle), coupledFoam time split into assembly, linear solve, turbulence and other (bottom).}{fig:hist:T2_np4:wall}
\cfhistfigure{hist_T2_np4_linear}{T2 backward-facing step, 4 ranks: linear solver and pseudo-time controls per outer iteration; bottom: simpleFoam linear iterations per field.}{fig:hist:T2_np4:linear}
\clearpage

\subsection{T3 airFoil2D, SST}\label{app:hist:T3_kOmegaSST_np1}

\begin{itemize}\setlength{\itemsep}{0pt}\small
\item coupledFoam: 3000 iterations, 683.5\,s wall-clock, 0.19 CPU-h; time per iteration 0.181\,s (first 10\,\%) and 0.142\,s (last 10\,\%), linear fit 0.234\,s/it ($R^2=0.999$).
\item simpleFoam: 20000 iterations, 586\,s wall-clock, 0.163 CPU-h; time per iteration 0.03\,s (first 10\,\%) and 0.03\,s (last 10\,\%), linear fit 0.029\,s/it ($R^2=1.000$).
\item coupledFoam linear iterations in total: 56679.
\item simpleFoam D-042 window $W=10000$: $\bar C_d=0.0907$, $\bar C_l=0.2527$, stationary from iteration 10100.
\item coupledFoam D-042 window $W=1500$: $\bar C_d=0.1473$, $\bar C_l=0.2841$, stationary from iteration 2400.
\end{itemize}

\cfhistfigure{hist_T3_kOmegaSST_np1_loads}{T3 airFoil2D, SST: monitored quantities per outer iteration (left) and over wall-clock time (right); thin: raw, dashed: running mean over the D-042 window $W$ with $\pm$RMS band, shaded: final averaging window, dash-dot: iterations to a stationary window.}{fig:hist:T3_kOmegaSST_np1:loads}
\cfhistfigure{hist_T3_kOmegaSST_np1_residuals}{T3 airFoil2D, SST: residuals per outer iteration (left) and over wall-clock time (right), coupledFoam top, simpleFoam bottom, same y range.}{fig:hist:T3_kOmegaSST_np1:residuals}
\cfhistfigure{hist_T3_kOmegaSST_np1_wall}{T3 airFoil2D, SST: cumulative wall time with linear fit (top), wall time per iteration with rolling median (middle), coupledFoam time split into assembly, linear solve, turbulence and other (bottom).}{fig:hist:T3_kOmegaSST_np1:wall}
\cfhistfigure{hist_T3_kOmegaSST_np1_linear}{T3 airFoil2D, SST: linear solver and pseudo-time controls per outer iteration; bottom: simpleFoam linear iterations per field.}{fig:hist:T3_kOmegaSST_np1:linear}
\clearpage

\subsection{T3 airFoil2D, GEKO}\label{app:hist:T3_GEKO_np1}

\begin{itemize}\setlength{\itemsep}{0pt}\small
\item coupledFoam: 3000 iterations, 821.5\,s wall-clock, 0.228 CPU-h; time per iteration 0.185\,s (first 10\,\%) and 0.192\,s (last 10\,\%), linear fit 0.27\,s/it ($R^2=0.999$).
\item simpleFoam: 20000 iterations, 602\,s wall-clock, 0.167 CPU-h; time per iteration 0.03\,s (first 10\,\%) and 0.03\,s (last 10\,\%), linear fit 0.03\,s/it ($R^2=1.000$).
\item coupledFoam linear iterations in total: 67769.
\item simpleFoam D-042 window $W=10000$: $\bar C_d=0.0343$, $\bar C_l=0.8378$, stationary from iteration 10000.
\item coupledFoam D-042 window $W=1500$: $\bar C_d=0.0581$, $\bar C_l=0.7509$, stationary from iteration 1650.
\end{itemize}

\cfhistfigure{hist_T3_GEKO_np1_loads}{T3 airFoil2D, GEKO: monitored quantities per outer iteration (left) and over wall-clock time (right); thin: raw, dashed: running mean over the D-042 window $W$ with $\pm$RMS band, shaded: final averaging window, dash-dot: iterations to a stationary window.}{fig:hist:T3_GEKO_np1:loads}
\cfhistfigure{hist_T3_GEKO_np1_residuals}{T3 airFoil2D, GEKO: residuals per outer iteration (left) and over wall-clock time (right), coupledFoam top, simpleFoam bottom, same y range.}{fig:hist:T3_GEKO_np1:residuals}
\cfhistfigure{hist_T3_GEKO_np1_wall}{T3 airFoil2D, GEKO: cumulative wall time with linear fit (top), wall time per iteration with rolling median (middle), coupledFoam time split into assembly, linear solve, turbulence and other (bottom).}{fig:hist:T3_GEKO_np1:wall}
\cfhistfigure{hist_T3_GEKO_np1_linear}{T3 airFoil2D, GEKO: linear solver and pseudo-time controls per outer iteration; bottom: simpleFoam linear iterations per field.}{fig:hist:T3_GEKO_np1:linear}
\clearpage

\subsection{T4a motorBike, 354k cells, 10 ranks}\label{app:hist:T4a_np10}

\begin{itemize}\setlength{\itemsep}{0pt}\small
\item coupledFoam: run in progress, skipped (pending).
\item simpleFoam: 3000 iterations, 1516\,s wall-clock, 4.21 CPU-h; time per iteration 0.52\,s (first 10\,\%) and 0.5\,s (last 10\,\%), linear fit 0.502\,s/it ($R^2=1.000$).
\item simpleFoam D-042 window $W=2250$: $\bar C_d=0.3964$, $\bar C_l=0.0768$, stationary from iteration 2300.
\item simpleFoam: the monitor has 4500 iterations, the solver log 3000; the later iterations have no wall-clock time (right column).
\end{itemize}

\cfhistfigure{hist_T4a_np10_loads}{T4a motorBike, 354k cells, 10 ranks: monitored quantities per outer iteration (left) and over wall-clock time (right); thin: raw, dashed: running mean over the D-042 window $W$ with $\pm$RMS band, shaded: final averaging window, dash-dot: iterations to a stationary window.}{fig:hist:T4a_np10:loads}
\cfhistfigure{hist_T4a_np10_residuals}{T4a motorBike, 354k cells, 10 ranks: residuals per outer iteration (left) and over wall-clock time (right), coupledFoam top, simpleFoam bottom, same y range.}{fig:hist:T4a_np10:residuals}
\cfhistfigure{hist_T4a_np10_wall}{T4a motorBike, 354k cells, 10 ranks: cumulative wall time with linear fit (top), wall time per iteration with rolling median (middle), coupledFoam time split into assembly, linear solve, turbulence and other (bottom).}{fig:hist:T4a_np10:wall}
\cfhistfigure{hist_T4a_np10_linear}{T4a motorBike, 354k cells, 10 ranks: linear solver and pseudo-time controls per outer iteration; bottom: simpleFoam linear iterations per field.}{fig:hist:T4a_np10:linear}
\clearpage

\subsection{T4b motorBike, 1.70M cells, 10 ranks}\label{app:hist:T4b_np10}

\begin{itemize}\setlength{\itemsep}{0pt}\small
\item coupledFoam: 64 iterations, n/a\,s wall-clock, n/a CPU-h; time per iteration 3.47\,s (first 10\,\%) and 5.14\,s (last 10\,\%), linear fit 4.94\,s/it ($R^2=0.999$).
\item simpleFoam: 4000 iterations, 1.188e+04\,s wall-clock, 33 CPU-h; time per iteration 2.87\,s (first 10\,\%) and 2.91\,s (last 10\,\%), linear fit 2.96\,s/it ($R^2=1.000$).
\item coupledFoam linear iterations in total: 284.
\item simpleFoam D-042 window $W=2000$: $\bar C_d=0.3996$, $\bar C_l=0.0657$, stationary from iteration 2050.
\item coupledFoam D-042 window $W=64$: $\bar C_d=0.4974$, $\bar C_l=0.0461$, stationary from iteration n/a.
\item constant over the whole coupledFoam run (panel omitted): cuts = 0, eta = 0.5, nRollback = 0.
\end{itemize}

\cfhistfigure{hist_T4b_np10_loads}{T4b motorBike, 1.70M cells, 10 ranks: monitored quantities per outer iteration (left) and over wall-clock time (right); thin: raw, dashed: running mean over the D-042 window $W$ with $\pm$RMS band, shaded: final averaging window, dash-dot: iterations to a stationary window.}{fig:hist:T4b_np10:loads}
\cfhistfigure{hist_T4b_np10_residuals}{T4b motorBike, 1.70M cells, 10 ranks: residuals per outer iteration (left) and over wall-clock time (right), coupledFoam top, simpleFoam bottom, same y range.}{fig:hist:T4b_np10:residuals}
\cfhistfigure{hist_T4b_np10_wall}{T4b motorBike, 1.70M cells, 10 ranks: cumulative wall time with linear fit (top), wall time per iteration with rolling median (middle), coupledFoam time split into assembly, linear solve, turbulence and other (bottom).}{fig:hist:T4b_np10:wall}
\cfhistfigure{hist_T4b_np10_linear}{T4b motorBike, 1.70M cells, 10 ranks: linear solver and pseudo-time controls per outer iteration; bottom: simpleFoam linear iterations per field.}{fig:hist:T4b_np10:linear}
\clearpage

\subsection{T5 Ahmed body, 10 ranks}\label{app:hist:T5_np10}

\begin{itemize}\setlength{\itemsep}{0pt}\small
\item coupledFoam: run not available yet (pending).
\item simpleFoam: run not available yet (pending).
\end{itemize}

\cfhistfigure{hist_T5_np10_loads}{T5 Ahmed body, 10 ranks: monitored quantities per outer iteration (left) and over wall-clock time (right); thin: raw, dashed: running mean over the D-042 window $W$ with $\pm$RMS band, shaded: final averaging window, dash-dot: iterations to a stationary window.}{fig:hist:T5_np10:loads}
\cfhistfigure{hist_T5_np10_residuals}{T5 Ahmed body, 10 ranks: residuals per outer iteration (left) and over wall-clock time (right), coupledFoam top, simpleFoam bottom, same y range.}{fig:hist:T5_np10:residuals}
\cfhistfigure{hist_T5_np10_wall}{T5 Ahmed body, 10 ranks: cumulative wall time with linear fit (top), wall time per iteration with rolling median (middle), coupledFoam time split into assembly, linear solve, turbulence and other (bottom).}{fig:hist:T5_np10:wall}
\cfhistfigure{hist_T5_np10_linear}{T5 Ahmed body, 10 ranks: linear solver and pseudo-time controls per outer iteration; bottom: simpleFoam linear iterations per field.}{fig:hist:T5_np10:linear}
\clearpage
}{%
  \pending{The appendix text \texttt{figures/histories.tex} is written by
  \texttt{bench/make\_report.py}.}}

% =============================================================================
\clearpage
\section{Missing or stale results}\label{app:missing}
% =============================================================================

Every number, table and figure of this document is generated by
\code{bench/make\_report.py} from results produced at one frozen commit of
the code. The generator refuses results of any other commit, results
written from a modified source tree and run directories without a
provenance record; they are shown as ``pending'' in the text and listed
here. The list is generated together with the rest of the data.

\IfFileExists{tables/missing_results.tex}{% Generated by bench/make_report.py - do not edit
Target commit \texttt{bd128a7} (\texttt{bd128a73f1339ab056571600b93fe640378e32ed}); items not produced at the target commit are left out and shown as ``pending'' in the text. A result counts when its record (\texttt{results/*/<name>.json}, field \texttt{gitCommit}) or its run directory (\texttt{provenance.json}, written by the test harness at the start of every run) names this commit and a clean source tree. The \code{simpleFoam} reference runs are exempt from the run check (untouched system OpenFOAM, cached). Items listed: 64.

{\small
\begin{longtable}{@{}p{0.21\linewidth}p{0.31\linewidth}p{0.15\linewidth}p{0.1\linewidth}p{0.1\linewidth}@{}}
\caption{Missing or stale results (generated).}\label{tab:missing}\\
\toprule case / record & item (reason) & found commit & date & status\\ \midrule\endfirsthead
\toprule case / record & item (reason) & found commit & date & status\\ \midrule\endhead
\bottomrule\endlastfoot
\texttt{T-fpe\_*} & result record \newline{\footnotesize (no result)} & {\footnotesize\texttt{none}} &  & pending\\
\texttt{T-restart\_T1} & result record \newline{\footnotesize (no result)} & {\footnotesize\texttt{none}} &  & pending\\
\texttt{T-restart\_T3-SST} & result record \newline{\footnotesize (no result)} & {\footnotesize\texttt{none}} &  & pending\\
\texttt{T0\_Re1000\_np1} & result record \newline{\footnotesize (written from a modified tree)} & {\footnotesize\texttt{9090388-dirty}} & 2026-09-22\newline 13:20 & pending\\
\texttt{T0\_Re1000\_np1} & run directory \newline{\footnotesize (run has no provenance.json)} & {\footnotesize\texttt{none}} & 2026-09-22\newline 13:20 & pending\\
\texttt{T0\_Re1000\_np4} & result record \newline{\footnotesize (record from another commit)} & {\footnotesize\texttt{93dc500}} & 2026-09-22\newline 09:24 & pending\\
\texttt{T0\_Re1000\_np4} & run directory \newline{\footnotesize (run has no provenance.json)} & {\footnotesize\texttt{none}} & 2026-09-22\newline 13:22 & pending\\
\texttt{T0\_Re100\_np1} & result record \newline{\footnotesize (record from another commit)} & {\footnotesize\texttt{cb5ea56}} & 2026-09-22\newline 11:53 & pending\\
\texttt{T0\_Re100\_np1} & run directory \newline{\footnotesize (run has no provenance.json)} & {\footnotesize\texttt{none}} & 2026-09-22\newline 13:19 & pending\\
\texttt{T0\_Re100\_np4} & result record \newline{\footnotesize (record from another commit)} & {\footnotesize\texttt{93dc500}} & 2026-09-22\newline 09:23 & pending\\
\texttt{T0\_Re100\_np4} & run directory \newline{\footnotesize (run has no provenance.json)} & {\footnotesize\texttt{none}} & 2026-09-22\newline 13:20 & pending\\
\texttt{T1\_np1} & result record \newline{\footnotesize (written from a modified tree)} & {\footnotesize\texttt{07abc64-dirty}} & 2026-09-22\newline 09:51 & pending\\
\texttt{T1\_np1} & run directory \newline{\footnotesize (run has no provenance.json)} & {\footnotesize\texttt{none}} & 2026-09-22\newline 13:22 & pending\\
\texttt{T1\_np4} & result record \newline{\footnotesize (written from a modified tree)} & {\footnotesize\texttt{07abc64-dirty}} & 2026-09-22\newline 09:51 & pending\\
\texttt{T1\_np4} & run directory \newline{\footnotesize (run has no provenance.json)} & {\footnotesize\texttt{none}} & 2026-09-22\newline 13:22 & pending\\
\texttt{T2\_np1} & result record \newline{\footnotesize (record from another commit)} & {\footnotesize\texttt{cce9ede}} & 2026-09-22\newline 09:55 & pending\\
\texttt{T2\_np1} & run directory \newline{\footnotesize (run has no provenance.json)} & {\footnotesize\texttt{none}} & 2026-09-22\newline 13:27 & pending\\
\texttt{T2\_np4} & result record \newline{\footnotesize (record from another commit)} & {\footnotesize\texttt{cce9ede}} & 2026-09-22\newline 09:56 & pending\\
\texttt{T2\_np4} & run directory \newline{\footnotesize (run has no provenance.json)} & {\footnotesize\texttt{none}} & 2026-09-22\newline 09:56 & pending\\
\texttt{T3\_GEKO\_np1} & result record \newline{\footnotesize (record from another commit)} & {\footnotesize\texttt{4e7e80d}} & 2026-09-22\newline 03:48 & pending\\
\texttt{T3\_GEKO\_np1} & run directory \newline{\footnotesize (run has no provenance.json)} & {\footnotesize\texttt{none}} & 2026-09-22\newline 03:48 & pending\\
\texttt{T3\_GEKO\_np4} & result record \newline{\footnotesize (no result)} & {\footnotesize\texttt{none}} &  & pending\\
\texttt{T3\_kOmegaSST\_np1} & result record \newline{\footnotesize (record from another commit)} & {\footnotesize\texttt{4e7e80d}} & 2026-09-22\newline 03:24 & pending\\
\texttt{T3\_kOmegaSST\_np1} & run directory \newline{\footnotesize (run has no provenance.json)} & {\footnotesize\texttt{none}} & 2026-09-22\newline 03:24 & pending\\
\texttt{T3\_kOmegaSST\_np4} & result record \newline{\footnotesize (no result)} & {\footnotesize\texttt{none}} &  & pending\\
\texttt{T4a\_np10} & result record \newline{\footnotesize (written from a modified tree)} & {\footnotesize\texttt{ca33547-dirty}} & 2026-09-22\newline 11:20 & pending\\
\texttt{T4a\_np10} & run directory \newline{\footnotesize (run has no provenance.json)} & {\footnotesize\texttt{none}} & 2026-09-22\newline 13:40 & pending\\
\texttt{T4b\_np10} & run directory \newline{\footnotesize (run has no provenance.json)} & {\footnotesize\texttt{none}} & 2026-09-22\newline 13:41 & pending\\
\texttt{T4b\_np<n>} & result record \newline{\footnotesize (no result)} & {\footnotesize\texttt{none}} &  & pending\\
\texttt{T5\_np<n>} & result record \newline{\footnotesize (no result)} & {\footnotesize\texttt{none}} &  & pending\\
\texttt{T\_scaling\_T4b} & result record \newline{\footnotesize (no result)} & {\footnotesize\texttt{none}} &  & pending\\
\texttt{Test-block4Ops} & result record \newline{\footnotesize (written from a modified tree)} & {\footnotesize\texttt{9f15f8f-dirty}} & 2026-09-22\newline 00:49 & pending\\
\texttt{Test-blockFGMRES} & result record \newline{\footnotesize (written from a modified tree)} & {\footnotesize\texttt{9f15f8f-dirty}} & 2026-09-22\newline 00:49 & pending\\
\texttt{Test-blockGAMG} & result record \newline{\footnotesize (written from a modified tree)} & {\footnotesize\texttt{9f15f8f-dirty}} & 2026-09-22\newline 00:49 & pending\\
\texttt{Test-blockGAMG\_cycles} & result record \newline{\footnotesize (written from a modified tree)} & {\footnotesize\texttt{9f15f8f-dirty}} & 2026-09-22\newline 00:49 & pending\\
\texttt{Test-blockGAMG\_np1} & result record \newline{\footnotesize (record carries no gitCommit)} & {\footnotesize\texttt{none}} &  & pending\\
\texttt{Test-blockGAMG\_np4} & result record \newline{\footnotesize (record carries no gitCommit)} & {\footnotesize\texttt{none}} &  & pending\\
\texttt{Test-blockMatrix} & result record \newline{\footnotesize (written from a modified tree)} & {\footnotesize\texttt{9f15f8f-dirty}} & 2026-09-22\newline 00:49 & pending\\
\texttt{Test-blockMatrix\_np1} & result record \newline{\footnotesize (record carries no gitCommit)} & {\footnotesize\texttt{none}} &  & pending\\
\texttt{Test-blockMatrix\_np4} & result record \newline{\footnotesize (record carries no gitCommit)} & {\footnotesize\texttt{none}} &  & pending\\
\texttt{Test-convergenceMonitor} & result record \newline{\footnotesize (record from another commit)} & {\footnotesize\texttt{cb5ea56}} & 2026-09-22\newline 11:53 & pending\\
\texttt{Test-doubleReduce} & result record \newline{\footnotesize (written from a modified tree)} & {\footnotesize\texttt{9f15f8f-dirty}} & 2026-09-22\newline 00:49 & pending\\
\texttt{Test-doubleReduce\_np1} & result record \newline{\footnotesize (record carries no gitCommit)} & {\footnotesize\texttt{none}} &  & pending\\
\texttt{Test-doubleReduce\_np4} & result record \newline{\footnotesize (record carries no gitCommit)} & {\footnotesize\texttt{none}} &  & pending\\
\texttt{Test-precision} & result record \newline{\footnotesize (record carries no gitCommit)} & {\footnotesize\texttt{none}} &  & pending\\
\texttt{benchmark A-H} & results/bench/*.json \newline{\footnotesize (no result)} & {\footnotesize\texttt{none}} &  & pending\\
\texttt{diagnostics\_force\_stats} & result record \newline{\footnotesize (record from another commit)} & {\footnotesize\texttt{cb5ea56}} & 2026-09-22\newline 11:53 & pending\\
\texttt{diagnostics\_level0} & result record \newline{\footnotesize (record from another commit)} & {\footnotesize\texttt{cb5ea56}} & 2026-09-22\newline 11:53 & pending\\
\texttt{diagnostics\_np1} & result record \newline{\footnotesize (record from another commit)} & {\footnotesize\texttt{cb5ea56}} & 2026-09-22\newline 11:53 & pending\\
\texttt{exploratory decisions} & results/exploratory/decisions.json \newline{\footnotesize (record carries no gitCommit)} & {\footnotesize\texttt{none}} & 2026-09-21\newline 23:13 & pending\\
\texttt{exploratory linsolver} & results/exploratory/linsolver.json \newline{\footnotesize (record carries no gitCommit)} & {\footnotesize\texttt{none}} & 2026-09-21\newline 23:13 & pending\\
\texttt{exploratory re1000} & results/exploratory/re1000.json \newline{\footnotesize (record carries no gitCommit)} & {\footnotesize\texttt{none}} & 2026-09-21\newline 23:13 & pending\\
\texttt{exploratory symbol} & results/exploratory/symbol.json \newline{\footnotesize (record carries no gitCommit)} & {\footnotesize\texttt{none}} & 2026-09-21\newline 23:13 & pending\\
\texttt{render\_T4a\_delta} & ParaView render (no render provenance) \newline{\footnotesize (record carries no gitCommit)} & {\footnotesize\texttt{none}} &  & pending\\
\texttt{render\_T4a\_geometry} & ParaView render (no render provenance) \newline{\footnotesize (record carries no gitCommit)} & {\footnotesize\texttt{none}} &  & pending\\
\texttt{render\_T4a\_mean\_U} & ParaView render (no render provenance) \newline{\footnotesize (record carries no gitCommit)} & {\footnotesize\texttt{none}} &  & pending\\
\texttt{render\_T4a\_slice\_U} & ParaView render (no render provenance) \newline{\footnotesize (record carries no gitCommit)} & {\footnotesize\texttt{none}} &  & pending\\
\texttt{render\_T4a\_slice\_p} & ParaView render (no render provenance) \newline{\footnotesize (record carries no gitCommit)} & {\footnotesize\texttt{none}} &  & pending\\
\texttt{render\_T4a\_surface\_Cp} & ParaView render (no render provenance) \newline{\footnotesize (record carries no gitCommit)} & {\footnotesize\texttt{none}} &  & pending\\
\texttt{render\_T4b\_geometry} & ParaView render (no render provenance) \newline{\footnotesize (record carries no gitCommit)} & {\footnotesize\texttt{none}} &  & pending\\
\texttt{render\_T4b\_slice\_U} & ParaView render (no render provenance) \newline{\footnotesize (record carries no gitCommit)} & {\footnotesize\texttt{none}} &  & pending\\
\texttt{render\_T4b\_slice\_p} & ParaView render (no render provenance) \newline{\footnotesize (record carries no gitCommit)} & {\footnotesize\texttt{none}} &  & pending\\
\texttt{render\_T4b\_surface\_Cp} & ParaView render (no render provenance) \newline{\footnotesize (record carries no gitCommit)} & {\footnotesize\texttt{none}} &  & pending\\
\texttt{test\_env} & result record \newline{\footnotesize (written from a modified tree)} & {\footnotesize\texttt{03954b0-dirty}} & 2026-09-21\newline 20:33 & pending\\
\end{longtable}
}
}{%
  \pending{The list \texttt{tables/missing\_results.tex} is written by
  \texttt{bench/make\_report.py}.}}

\clearpage
\bibliographystyle{plainnat}
\bibliography{references}

\end{document}
